\chapter{General Mass Law: absolute action scale \texorpdfstring{\(10^{-34}\)}{10⁻³⁴}, domain \texorpdfstring{\(81/56/13/324\)}{81/56/13/324}, null sector and operatorial construction of the material spectrum}
\label{chap:ley-general-masas-dominio-completo}

\section{Operator Domain of the General Mass Law: 81 Cells, 56 Families, 13 Multisections, 324 Routes, and Null Sector}
The General Mass Law acts on persistent sections of the oriented atlas, not on a table of species. Its internal domain contains \(81\) cells, \(56\) route families, \(13\) multisections, and \(324\) oriented routes, in addition to a separate null chart. The censuses of physical outputs are subsequent and do not alter this domain.

\section{Determinantal origin of the absolute scale}
The integral form \(H_4\) has Smith normal form
\(\operatorname{diag}(1,2,2,4)\), so its cokernel has order \(16\).
It follows that
\[
 \Sigma_H=\varphi_{\rm HMT}\alpha_{\rm HMT}^{16},
 \qquad \operatorname{ord}_{10}(\Sigma_H)=-34.
\]
The absolute chain is
\[
 \Sigma_H\to\hbar_{\rm pre/ret}\to\omega_0
 \to\varepsilon_{\rm clk}\to m_{\rm clk}
 \to\widehat M^{(0)}\to\widehat M^{(\beta)}.
\]
The decade \(-34\) is never introduced from the SI system or absorbed into the
relative correction.

\section{Governing construction of the law and the material domain}
The full development of the law is organized here by result: state, route,
record, signature, character, towers, absolute branch, and operator. Its internal order
is preserved, while the realizations by species are reserved for the corresponding
chapter.

\begingroup
\let\chapter\section
\let\section\subsection
\let\subsection\subsubsection
\let\subsubsection\paragraph
\let\clearpage\relax
\section{Spectral Definition of Mass as a Persistence Measure of a Route \(1\) Coupled to a Physical Regime}
Mass is the spectral measure of the persistence of a route when a genealogical section is coupled to a physical regime. It is not a scalar attribute affixed in advance to a particle name. The APP--TRIT--TPK dynamics constructs the route, its memory, its signature, and its character; Dimensional Mechanics realizes that genealogy on a resonance mass line and determines, through the coupling, the observable subspace and its spectrum.

An elementary particle is represented by a persistent section; a noncentral family, by a stable multisection; a composite particle, by a composition of genealogical registers with a linking boundary; a resonance, by a pole of the coupled operator; a topological sector, by its fusion ring and its braiding representation before any mass chart; a virtual section, by its terminal horizon of prolongation. These realizations are distinct and must not be flattened into a single list of numbers.

The external magnitude never enters the selection of cell, route, signature, reader, tower coefficient, or eigenvector. Metrological comparison is performed only after the internal output has been constructed and fixed.

\Needspace{7\baselineskip}
\section{Complete state space of mass}
Let \(\mathcal C_{81}\) be the cellular atlas. Each cell admits four bidirectional route orientations, so the finite space of oriented routes is

\[
 \mathcal R_{324}
 =\{(c,h_0,v_0):c\in\mathcal C_{81},\ (h_0,v_0)\in\{\pm1\}^2\}.
\]
The cellular projection, the family relation, and the coarse stratification produce

\[
 |\mathcal C_{81}|=81,\qquad
 |\mathcal F_{56}|=56,\qquad
 |\mathcal M_{13}|=13,\qquad
 |\mathcal R_{324}|=324.
\]
The electron occupies the unique central fixed point. The twelve noncentral multisections are finite fibers stable under cellular inversion and do not, in general, admit an equivariant point section. Their natural object is therefore an operator on the fiber rather than a number chosen from among its cells. The null sector is incorporated through the zero section of the mass line, not as the value of a positive character.

The physical output catalog distinguishes three charged leptons, three neutrino interaction realizations, six quark flavors with their color orbits, two null gauge realizations, two massive gauge bosons, one scalar sector, two composite families, three topological codomains, and a system of virtual sections. This subsequent classification does not limit the number of composite or resonant states: composition and excitation generate additional families within their respective codomains.

\Needspace{7\baselineskip}
\section{Genealogical trajectory of a mass realization}
The law does not begin with a particle label nor with a number. Its order is

\[
\boxed{
\begin{aligned}
\mathrm{APP}\to\mathrm{TRIT}\to\mathrm{TPK}
&\to\widetilde\Gamma^{\rm surv}
\to\Gamma^{\rm obs}
\to L_{108}
\to\sigma\in\mathbb Z^5\\
&\to\chi_0
\to\chi_{\rm full}
\to\widehat{\mathcal M}
\to\text{coupling}\\
&\to\text{spectrum/pole}
\to\mathcal L_M
\to\text{comparison}.
\end{aligned}}
\]
Here:

\begin{itemize}
\item APP arranges the local incompatibility, the sum/product sheets,
and their orientation;
\item the TRIT preserves sign, threshold, and bidirectional opening;
\item the TPK evolves phase, memory, carry, survival, and return;
\item a particle is a persistent section of an extension, and its observable route is
the history it leaves in the genealogical record;

\item the integer signature is additive;
\item the character turns additive composition into a positive ratio;
\item the towers refine the character through a single global harmonic operator;
\item the mass line supplies the dimensional type;
\item measurement is coupling: the coupling determines which subspace and which spectrum are
physically accessible, without consulting the quantity to be measured subsequently.

\end{itemize}
Mass does not select the route. The persistent route, the genealogical record, and the coupling select the spectral object whose mass is read.

\Needspace{7\baselineskip}
\subsection{Conserved components of the enriched state of the TPK along the mass trajectory}
The state that feeds the mass law is

\[
 X_K=(w_6,w_{30},R_{36},g_K,B_K,p_K,c_K,\varepsilon_K,\omega_K),
\]
where short word, prolonged word, region, nonadic phase, block memory, prefix, carry, sheet, and winding are preserved. The chain

\[
 w_6\longrightarrow w_{30}\longrightarrow R_{36}
 \longrightarrow G_9
\]
It does not directly project a mass: it constructs the persistent identity that is later read.

\Needspace{7\baselineskip}
\subsection{Nonadic filtration with memory}
From each state one obtains a refinement fiber formed by \(729\) subcylinders,

\[
 \{C_{K,j}:0\le j<729\}.
\]
The pruning \(G_9\) satisfies

\[
 G_9(C_{K,j})\in\{0,1\},\qquad
 \sum_{j=0}^{728}G_9(C_{K,j})=1.
\]
If \(j_K\) is the survivor,

\[
 P_{K+1}=729P_K+j_K.
\]
After nine steps,

\[
 g(K+9)=g(K),
\]
but they change, in general,

\[
 (B_{K+9},P_{K+9},\operatorname{pref}_{K+9},
 \operatorname{Witt}_{K+9})
 \ne
 (B_K,P_K,\operatorname{pref}_K,\operatorname{Witt}_K).
\]
This is the equational form of persistence: the phase returns, not the state. The mass can recognize a particle after the return because the genealogical record preserves the difference.

\Needspace{7\baselineskip}
\subsection{Nonadic carry transduction to the signed dodecaphasic state}
For a temporal coordinate \(\theta\), define

\[
 g_0(\theta)=\theta\bmod9,\qquad
 \beta(\theta)=
 \left\lfloor\frac{\widetilde\theta}{9}\right\rfloor\bmod12.
\]
Then

\[
 g_0(\theta+9)=g_0(\theta),\qquad
 \beta(\theta+9)=\beta(\theta)+1\pmod{12}.
\]
Thus, the nonadic return preserves the phase and advances one memory tooth. The 108 times are

\[
 108=9\cdot12,
\]
so that each genealogical record contains a complete cycle of both scales. The 36.864 carry compositions describe the local combinations of this bidirectional memory.

\Needspace{7\baselineskip}
\subsection{Completion and Conservation of the Unit}
The base boundary mil decomposes as

\[
 999=729+270=729+243+27,
\]
while the completion is

\[
 1000=729+271=729+243+27+1,
\qquad
 271=270+1.
\]
\(+1\) is the unique vertex lying entirely on the boundary and realizes the carry that returns the unit. The partition \(729/271\) distinguishes the active kernel from the mediating boundary; it is not a golden-ratio approximation. The same conservation allows a route to change blocks without losing its normalized identity.

\Needspace{7\baselineskip}
\subsection{Theorem of the complete oriented atlas}
Every cell \(c\in\mathcal C_{81}\) has four orientations

\[
 (h_0,v_0)\in\{\pm1\}^2.
\]
The map

\[
 (c,h_0,v_0)\longmapsto\Gamma(c,h_0,v_0)
\]
is bijective over \(\mathcal R_{324}\). Consequently,

\[
 |\mathcal R_{324}|=4|\mathcal C_{81}|=324.
\]
Each route has 108 genealogical register entries, so the total temporal census is

\[
 81\cdot108=8748.
\]
The quadruple

\[
 (\text{seed row},\text{seed column},h_0,v_0)
\]
determines of manner single the route, its word \(w_6\), its word dodecaphasic, their returns, memory, carry, signature historical and readers. Therefore, the 324 routes not only may is enumerated: may is reconstructed and is joined exactly, field by field, without using a mass.

\Needspace{7\baselineskip}
\subsection{Quotients of families and multisections}
The quotient of the 81 cells by route history yields 56 families. Their multiplicities are distributed as

\[
 41\times1,\quad10\times2,\quad2\times3,\quad1\times4,\quad2\times5,
\]
and

\[
 41+20+6+4+10=81.
\]
The stratification into thirteen multisects possesses ranges

\[
 1,\quad2,\quad4,\quad6,\quad8,\quad12,\quad16
\]
with multiplicities

\[
 1,\quad2,\quad3,\quad2,\quad3,\quad1,\quad1,
\]
and sum of ranks

\[
 1+4+12+12+24+12+16=81.
\]
The electron occupies the rank-one block. The remaining sectors lie in higher-rank blocks and require a representation projector, not a nominal cell.

\Needspace{7\baselineskip}
\subsection{Signature, readers, and memory by route}
For each \(\Gamma\in\mathcal R_{324}\), the following are calculated:

\[
 \sigma_\Gamma\in\mathbb Z^5,\qquad
 K_D(\Gamma),\qquad K_\Omega(\Gamma),
\]
together with first returns of class, orbit, cell, and kinematic state; separation of prefixes; carry debt; ambiguity; stabilized triads; sheet and winding. The family \(K_D\) has 14 profiles; \(K_\Omega\), nine; together they form 40 pairs and 18 coincidences. This diversity is the fine information that makes it possible to pass from a coarse fiber to a physical operator.

\Needspace{7\baselineskip}
\section{From the APP defect to the full signature}
\Needspace{7\baselineskip}
\subsection{Primordial defect and unfolding}
At the center APP,

\[
 B_c=9P_++5P_-,\qquad 9-5=4.
\]
The primordial local gap is 4. The global defect \(54\) belongs to a later unfolding \(2\to6\to54\); it must not be projected back onto the center as though it were the same operation.

\Needspace{7\baselineskip}
\subsection{\(12\) eventos and divisor dodecafasico}
The oriented coordinate does not receive a factor \(1/6\) arbitrarily. The chain is

\[
12\ \text{events}
\to s_C^{\rm tot}
\to 1+5+6
\to s_C^{(12)}=\frac{s_C^{\rm tot}}6
\to W_\pm=6n_A\pm s_C
\to\operatorname{SNF}(1,12)
\to\chi_{\rm mass}.
\]
The \(1/6\) is the oriented quotient of the closure of twelve events. It is not the residue (-1/12), it is not the pseudoinverse of a valency, and it is not the hexad--octad incidence.

\Needspace{7\baselineskip}
\subsection{Construction of the complete full signature from the route register}
The genealogical record CT--108 produces

\[
 \sigma(\Gamma)
 =\bigl(n_A,s_C,k_{\Delta_4},b_{120},b_\zeta\bigr)\in\mathbb Z^5.
\]
The coordinates have distinct origins:

\begin{itemize}
\item \(n_A\): dominant duration/persistence;
\item \(s_C\): net orientation of the bilayer, divided by six after the
dodecaphasic closure;

\item \(k_{\Delta_4}\): minimal-torsion channel;
\item \(b_{120}\): event coefficient of the register;
\item \(b_\zeta\): step-three autocorrelation.
\end{itemize}
We define

\[
 \zeta_{270}:=135\Delta_4^2,
\]
and remains separated from

\[
 s_{270}=q_-^{270}-q_+^{270}.
\]
The 324 oriented routes are joined bijectively to their five historical coordinates through the canonical route identity. This join uses only genealogical structure and consults no metrological magnitude.

\Needspace{7\baselineskip}
\subsection{Direct angle, counter-angle, and conjugate channels}
The direct angular coordinate proceeds from the already constructed fine structure constant:

\[
 A=1000\alpha_{\rm HMT}
 =7.297352569283800997285105472380663\ldots{}^\circ.
\]
The counter-angle is not obtained from a mass or from the Barbero--Immirzi functional. Let

\[
 D_A=\exp\!\left(-\frac{100\pi}{9}\alpha_{\rm HMT}\right),
\]
\[
 \mathcal C_\pi(\alpha_{\rm HMT})
 =169\alpha_{\rm HMT}^{6}
 +\frac{D_A\alpha_{\rm HMT}^{7}}
        {1-D_A\alpha_{\rm HMT}},
\]
\[
 y_*=\frac{\pi}{90}(H_5+\alpha_{\rm HMT})
      -\mathcal C_\pi(\alpha_{\rm HMT}),
 \qquad
 C^*=\frac{180}{\pi}y_*.
\]
Therefore,

\[
 C^*=2.123738338968645724261951380393\ldots{}^\circ.
\]
The coordinate \(A\) is the common angular mode and \(C^*\), the oriented differential mode. Its passage to the Archimedean chart is

\[
 x=A_{\rm rad}=\frac{\pi A}{180},
 \qquad
 y=C^*_{\rm rad}=\frac{\pi C^*}{180}.
\]
In the real algebra generated by the involution \(R=\operatorname{diag}(1,-1)\), with \(P_\pm=(I\pm R)/2\), the biangular block is

\[
 B_{A,C^*}=AI+C^*R
 =(A+C^*)P_++(A-C^*)P_-.
\]
Angle plus counter-angle and angle minus counter-angle are thus the two eigenvalues of a single object. Exponentiation produces the channels

\[
 q_+=e^{-x-y},\qquad q_-=e^{-x+y}.
\]
The chart is invertible:

\[
 x=-\frac12\log(q_+q_-),
 \qquad
 y=\frac12\log\frac{q_-}{q_+}.
\]
The transformation \(C^*\mapsto-C^*\) preserves the projectors and exchanges the channels. The return coordinate satisfies, after the regional construction,

\[
 \eta_{\rm ret}=\frac{C^*}{2}-\alpha_{\rm HMT}\,{\rm deg},
\]
without constituting a second selection rule for \(C^*\).

\Needspace{7\baselineskip}
\subsection{Central Character}
With

\[
 \Theta_0=
 \left(A_{\rm rad},\frac{C^*_{\rm rad}}6,
 \Delta_4,s_{120},\zeta_{270}\right),
\]
the central character is

\[
 \chi_0(\sigma)=
 \exp\!\left(
 n_AA_{\rm rad}+\frac{s_C}{6}C^*_{\rm rad}
 +k_{\Delta_4}\Delta_4+b_{120}s_{120}+b_\zeta\zeta_{270}
 \right).
\]
By additivity of the signature,

\[
 \chi_0(\sigma_1+\sigma_2)=\chi_0(\sigma_1)\chi_0(\sigma_2).
\]
The physical unit does not enter this exponential: \(\chi_0\) is a positive dimensionless automorphism of the mass line.

\Needspace{7\baselineskip}
\section{Global harmonic hierarchy of towers}
With the coordinates (x=\allowbreak{}A\_\{\textbackslash{}rm rad\}) and (y=\allowbreak{}C\textasciicircum{}*\_\{\textbackslash{}rm rad\}) already constructed, let

\[
 q_\pm=\exp\!\left[-\frac{\pi}{180}(A\pm C^*)\right]
 =e^{-x\mp y}.
\]
In a globally consistent convention,

\[
 T=e^{-xI-yR},\qquad
 c_n=\operatorname{Tr}(T^n)=q_-^n+q_+^n,
\]
\[
 s_n=-\operatorname{Tr}(RT^n)=q_-^n-q_+^n.
\]
Equivalently and operatorially,

\[
 c_n=2e^{-nx}\cosh(ny),\qquad
 s_n=2e^{-nx}\sinh(ny).
\]
The involution \(C^*\mapsto-C^*\) fixes the projectors and permutes the spectral coefficients, or equivalently the associated channels. It follows that

\[
 c_{m+n}=\frac{c_mc_n+s_ms_n}{2},\qquad
 s_{m+n}=\frac{s_mc_n+c_ms_n}{2},
\]
\[
 c_n^2-s_n^2=4e^{-2nx},
\]
and the common second-order recurrence. The admissible heights form

\[
 \mathfrak S=\langle90,120\rangle
 =\{90,120\}\cup\{30r:r\ge6\}.
\]
The values \(90,120,180,210,240,270,360,420\) are genealogical witnesses, not the complete tower. Before the quotient, height 360 preserves two histories:

\[
 k[\mathfrak S_0]
 \simeq k[X_{90},Y_{120}]/(X_{90}^4-Y_{120}^3).
\]
The general elevation is written

\[
 \chi_{\rm full}(\sigma,\eta)
 =\chi_0(\sigma)
  \exp\!\left(\sum_{h\in\mathfrak S}\eta_hs_h\right),
\]
with

\[
 N_{120}=b_{120}+\eta_{120}.
\]
The numerical sum does not erase the dual provenance: \(b_{120}\) comes from the ledger and \(\eta_{120}\) from the spectral refinement. Nor can it absorb \(b_\zeta\zeta_{270}\) into \(\eta_{270}s_{270}\).

\Needspace{7\baselineskip}
\section{Absolute Rama: action, clock and mass line}
The determinantal reader produces

\[
 \operatorname{SNF}(H_4)=\operatorname{diag}(1,2,2,4),\qquad
 |\operatorname{coker}H_4|=16,
\]
\[
 \Sigma_H=\varphi\alpha_{\rm HMT}^{16},\qquad
 \operatorname{ord}_{10}(\Sigma_H)=-34.
\]
The two oriented sections of a single action line are

\[
 \hbar_{\rm pre}=H_5\,10^{-34}\mathcal U_S,
 \qquad
 \hbar_{\rm ret}=\eta_{\rm ret}\,10^{-34}\mathcal U_S.
\]
Its additive fine chart and its multiplicative chart are

\[
 \delta_{2w}=H_5-\eta_{\rm ret},\qquad
 R_{\rm act}=\frac{H_5}{\eta_{\rm ret}}>0.
\]
They are not two Planck constants. The common decade fixes the absolute order; the quotient transports relative outward/return information.

Declared \(t_0\in\mathcal L_T^+\), CT--108 fixes

\[
 \omega_0=\frac{2\pi}{108t_0},\qquad
 \varepsilon_{\rm clk}=\hbar_{\rm ret}\omega_0,
 \qquad
 m_{\rm clk}=\varepsilon_{\rm clk}c_{\rm int}^{-2}.
\]
The character acts by homothety on \(\mathcal L_M\); it does not become \(\hbar\). In mass ratios, every common factor cancels:

\[
 \frac{m_Y}{m_X}
 =\chi_{\rm full}(\sigma_Y-\sigma_X,\eta_Y-\eta_X)
 R_{\rm act}^{\kappa_Y-\kappa_X}.
\]
Thus \(10^{-34}\) fixes the absolute scale but does not correct a relative residual. A relative modification can follow only from the route, the specific clock, the reader, or the towers.

Remains a distinction necessary. HMT constructs a section internal a time declared \(t_0\); the coordinate MeV requires a chart. For close the chart absolute without a convention externa must construirse

\[
 \Gamma\longmapsto t_\Gamma
 \quad\text{or}\quad
 \Gamma\longmapsto\omega_\Gamma.
\]
Section 10 defines a clock reader based on returns and holonomy. Its absolute normalization is determined when the route operator fixes a positive section of the temporal line.

\Needspace{7\baselineskip}

\endgroup

The normal form, the cokernel of order sixteen, the ablations, and the
distinction between \(h\) and \(\hbar\) are proved in full in
\cref{chap:escala-accion}. Here that already-constructed result is used to
fix the absolute branch of the law, without restarting its derivation.
\begingroup
\let\chapter\section
\let\section\subsection
\let\subsection\subsubsection
\let\subsubsection\paragraph
\let\clearpage\relax
\chapter{Conjugate angular coordinates and spectral decomposition}
\label{chap:canales}

The passage from the arithmetic closure of \(\alpha\) to the dimensional
hierarchy is effected through two angular coordinates. The first follows
directly from \(\alphaHMT\); the second is the oriented phase produced by
the regional analytic realization of the quintuple microfiber. Both reach
Dimensional Mechanics already constructed, before the exponential chart
that will generate the global moments is introduced.

\section{Transfer of the angular coordinates derived in HMT}

Parts I--III constructed distinct dependencies with a common root in
\(\alpha_{\rm HMT}\). The pre-return character follows from
\[
(\alpha_{\rm HMT},\varphi;\,9,5,4,12,54,\text{ graded rule})
\longrightarrow H_5,
\]
while the angular coordinate and the determinant follow from
\[
\alpha_{\rm HMT}\longrightarrow
A=1000\alpha_{\rm HMT}\,{}^\circ,
\qquad
(\alpha_{\rm HMT},\pi)\longrightarrow
D_A=e^{-100\pi\alpha_{\rm HMT}/9}.
\]
Through a distinct dependence, the regional microfiber fixes
\(\mathcal F_\pi\to\rho_\pi=169\), and with it
\[
\mathcal C_\pi(\alpha_{\rm HMT})
=169\alpha_{\rm HMT}^{6}
+\frac{D_A\alpha_{\rm HMT}^{7}}
       {1-D_A\alpha_{\rm HMT}}.
\]
The branch of \(H_5\) and the regional correction converge only at the
following step:
\[
\begin{aligned}
y_*&=\frac{\pi}{90}(H_5+\alpha_{\rm HMT})
     -\mathcal C_\pi(\alpha_{\rm HMT}),\\
C^*&=\frac{180}{\pi}y_*.
\end{aligned}
\]
Henceforth, \(A\) denotes the direct angular coordinate and \(C^*\) the
conjugate regional angular coordinate; these designations fix their types
without introducing a second selection rule.
Consequently, \(H_5\) is not an output of \(A\). Here \(y_*\) is the analytic
coordinate in radians and \(C^*\) its sexagesimal coordinate. Dimensional
Mechanics receives these already determined outputs; it does not redefine them.
The coordinate \(A\) is not a direct input to the formula of \(C^*\): the two
objects are coupled for the first time in \((A,C^*)\mapsto q_\pm\). The equality
\(\alpha_{\rm HMT}=A/1000\) is solely the inverse change of coordinate.
In this Part we retain the notations
\[
\begin{aligned}
 A&=7.297352569283800997285105472380663^\circ,\\
 \Cstar&=2.123738338968645724261951380393\ldots{}^\circ.
\end{aligned}
\]
The rounding used in the tables is
\(2.123738338968646^\circ\). The integer \(169\), the determinant \(D_A\), and
the regional recurrence belong to the preceding proof; neither any
SI comparison nor the hexad--octad functional intervenes in their generation.

The
angular identity
\[
 \etaRet=\frac{\Cstar}{2}-\alphaAngle,
 \qquad
 \alphaAngle:=\alphaHMT\,{\rm deg},
\]
defines the coordinate angular adimensional posterior to the return; does not  employed
for generate \(\Cstar\). Equivalently,
a time performed the construction regional,
\[
 \Cstar=2(\etaRet+\alphaAngle).
\]
This last equality is the inverse form of the same relation and not a
second selection rule for \(C^*\).

The regional realization corrects \(H_5\) before producing \(\etaRet\); therefore,
\(H_5\), \(\etaRet\), and \(\Cstar\) have distinct mathematical types. The symbol
\(\hbarHMT\) is reserved for the magnitude on the action line obtained only
after applying the scale \(10^{-34}\mathcal U_S\). Comparison with an SI mantissa belongs
to external metrological control and leaves the preceding construction invariant.

\section{Internal phases and exponential chart}

The abstract phases associated with the two angles are
\[
 a=\frac A{360},\qquad c=\frac{\Cstar}{360}
 \quad\text{in }\RR/\ZZ.
\]
Before applying a real exponential, the canonical representatives
\(\widetilde a,\widetilde c\in[0,1)\) are chosen. The exponential acts on those
representatives, not directly on the quotient classes.
Conversion to radians is introduced only upon choosing the
Archimedean chart
\[
 x:=A_{\rm rad}=\frac{\pi A}{180},\qquad
 y:=\Cstarrad=\frac{\pi\Cstar}{180},
\]
\[
 q_-=e^{-x+y},\qquad q_+=e^{-x-y}.
\]
The factor \(\pi/180\) belongs to this change analytic of coordinates. The
closure aritmetico of \(\alpha\), formulated before of this conversion, is
independent of that choice of chart angular.

\section{Central involutive algebra}

Consider the two-dimensional regular representation of the real algebra
generated by an involution. In a spectral basis, fix
\[
 R=\operatorname{diag}(1,-1),\qquad I=I_2,
\]
so that both eigenspaces have multiplicity one. Let
\[
 P_+=\frac{I+R}{2},\qquad P_-=\frac{I-R}{2}
\]
their spectral projectors. The central element determined by the
angular coordinates is
\[
 B_{A,\Cstar}=AI+\Cstar R
 =(A+\Cstar)P_++(A-\Cstar)P_-.
\]
The involution \(\Cstar\mapsto-\Cstar\) preserves the fixed projectors
\(P_\pm\) and interchanges the two spectral coefficients ---equivalently,
the associated channels---. Applying the exponential yields
\[
 T=e^{-xI-yR}=q_+P_++q_-P_-.
\]
Therefore, the two channels are the eigenvalues of a single element of the
algebra generated by \(I\) and \(R\).

\section{Invertibility of the channel chart}

\begin{theorem}[Reconstruction of the angular coordinates]
Para \(q_\pm>0\),
\[
 x=-\frac12\log(q_+q_-),
 \qquad
 y=\frac12\log\frac{q_-}{q_+}.
\]
Consequently, the map
\((A,\Cstar)\mapsto(q_+,q_-)\) is injective. The exchange
\(q_+\leftrightarrow q_-\) preserves \(A\) and transforms
\(\Cstar\mapsto-\Cstar\).
\end{theorem}
\begin{proof}
The multiplication and quotient of the two exponentials give
\(q_+q_-=e^{-2x}\) and \(q_-/q_+=e^{2y}\). Taking logarithms recovers
\(x\) and \(y\). Since \(\pi/180\neq0\), both angles are determined.
\end{proof}

\section{Addition and Recurrence Identities}

In this two-dimensional representation, the even and odd moments of the operator
\(T\) are
\[
 c_n=\Tr(T^n)=q_-^n+q_+^n,
 \qquad
 s_n=-\Tr(RT^n)=q_-^n-q_+^n.
\]

The multiplicity-one hypothesis is part of the domain of these formulas.
In a representation with multiplicities \(m_+\) and \(m_-\), the corresponding
traces would be
\(m_+q_+^n+m_-q_-^n\) and
\(m_-q_-^n-m_+q_+^n\), respectively.

\begin{theorem}[Laws of Composition of Moments]
For \(m,n\geq0\),
\[
 c_{m+n}=\frac{c_mc_n+s_ms_n}{2},
 \qquad
 s_{m+n}=\frac{s_mc_n+c_ms_n}{2},
\]
\[
 c_n^2-s_n^2=4e^{-2nx},
\]
and both sequences satisfy
\[
 X_{n+2}=c_1X_{n+1}-e^{-2x}X_n.
\]
Moreover,
\[
 \partial_Ac_n=-\frac{n\pi}{180}c_n,
 \quad
 \partial_As_n=-\frac{n\pi}{180}s_n,
\]
\[
 \partial_{\Cstar}c_n=\frac{n\pi}{180}s_n,
 \quad
 \partial_{\Cstar}s_n=\frac{n\pi}{180}c_n.
\]
\end{theorem}
\begin{proof}
Expand the products of
\(q_-^m\pm q_+^m\) with \(q_-^n\pm q_+^n\). The quadratic identity follows
from \((u+v)^2-(u-v)^2=4uv\). The recurrence corresponds to the characteristic
polynomial \(z^2-c_1z+q_-q_+\), and the derivatives are obtained using the
chain rule.
\end{proof}

\begin{corollary}[Determination of the complete hierarchy]
With \(A\) and \(\Cstar\) fixed, all pairs \((c_n,s_n)\) are
determined by the recurrence. No physical species introduces an additional
channel.
\end{corollary}

Finally,
\[
 s_n=2e^{-nx}\sinh(ny),
 \qquad
 c_n=2e^{-nx}\cosh(ny).
\]
The coordinate \(A\) determines the common decay, and \(\Cstar\) the oriented
separation. The decimal order of the action line is not introduced into these
formulas. It is deduced in \cref{chap:escala-accion} from the integral cokernel
of the local Hadamard block. The subsequent choice of a basis
expressed in \(\mathrm{J\,s}\) belongs to the metrological transducer and does not
alter the dimensionless identities.

\endgroup
\begingroup
\let\chapter\section
\let\section\subsection
\let\subsection\subsubsection
\let\subsubsection\paragraph
\let\clearpage\relax
\section{Action ellipse, radian, and double holonomy return}
\label{sec:elipse-accion-holonomia-rev8}
\index{action!ellipse}
\index{radian!action sector}
\index{Möbius!double return}

The preceding focal theorem fixes the angular form and the determinantal reader
fixes the action scale. Let us define
\[
 \eta=\operatorname{artanh}\frac{C^*}{A},
 \qquad
 a_S=\sqrt{2\hbar_{\rm HMT}^{\rm ret}e^\eta},
 \qquad
 b_S=\sqrt{2\hbar_{\rm HMT}^{\rm ret}e^{-\eta}}.
\]
The symplectic orbit
\[
 Q(\theta)=a_S\cos\theta,
 \qquad
 P(\theta)=b_S\sin\theta
\]
simultaneously preserves the shape ratio of the angular block and the area
fixed by the action line.

\begin{theorem}[Phase swept area]
\label{thm:area-fase-elipse-rev8}
For every \(\theta\),
\[
 \operatorname{Area}(\theta)
 =\frac12\int_0^\theta(QP'-PQ')\,dt
 =\hbar_{\rm HMT}^{\rm ret}\theta.
\]
Therefore,
\[
 \operatorname{Area}(1\ {\rm rad})=\hbar_{\rm HMT}^{\rm ret},
 \qquad
 \operatorname{Area}(2\pi)=h_{\rm HMT}^{\rm ret}.
\]
\end{theorem}

\begin{proof}
Since \(QP'-PQ'=a_Sb_S\) and
\(a_Sb_S=2\hbar_{\rm HMT}^{\rm ret}\), the integral equals
\(\hbar_{\rm HMT}^{\rm ret}\theta\). The second identity uses
\(h_{\rm HMT}^{\rm ret}=2\pi\hbar_{\rm HMT}^{\rm ret}\).
\end{proof}

APP supplies the arean oriented; \(A,C^*\) fix the form and the foci; the branch
determinantal fixes the scale. The turn primitive encloses \(h\) and a radian
of phase sweeps \(\hbar\). The radian ordinary remains
\(1\,\mathrm{rad}=180^\circ/\pi\); the construction adds its genealogy as
sector unitary of action, not another constant angular.

For a persistent section with holonomy
\(\varepsilon_\gamma\in\{+1,-1\}\), the condition
\[
 \varepsilon_\gamma e^{iJ_\gamma/\hbar}=1
\]
distinguishes two closures. In the cylinder, \(\varepsilon_\gamma=+1\), and the primitive
cycle transports \(J=2\pi\hbar=h\). In the Möbius band,
\(\varepsilon_\gamma=-1\): one visible turn changes sheet and transports
\(J=\pi\hbar=h/2\); the second restores orientation and completes \(h\). Per
radian of the internal phase, \(\hbar\) is transported in both cases; the factor
\(1/2\) belongs to the projection onto the visible turn of the Möbius base.

\begin{figure}[htbp]
\centering
\resizebox{.98\textwidth}{!}{\begin{tikzpicture}[
  font=\scriptsize,
  box/.style={draw,rounded corners=2mm,align=center,minimum height=.9cm,
    inner xsep=5pt},
  arr/.style={-{Stealth[length=2.2mm]},thick}]

  \begin{scope}[shift={(0,0)}]
    \draw[->,hmtgray] (-3.0,0)--(3.0,0) node[right]{\(u_+\)};
    \draw[->,hmtgray] (0,-2.0)--(0,2.0) node[above]{\(u_-\)};
    \draw[hmtblue,very thick] (0,0) ellipse (2.65cm and 1.45cm);
    \fill[hmtgold] (-1.85,0) circle (2pt) node[below]{\(F_-^{(\lambda)}\)};
    \fill[hmtgold] (1.85,0) circle (2pt) node[below]{\(F_+^{(\lambda)}\)};
    \node[text=hmtblue,font=\bfseries] at (0,2.55) {yield ellipse};
    \node[align=center,text=hmtgray] at (0,-2.4)
      {\(q_+=e^{-a_\lambda^2}\), \(q_-=e^{-b_\lambda^2}\)};
  \end{scope}

  \draw[arr,hmtgold] (3.3,0)--node[above,align=center]{same form\\action scale} (6.0,0);

  \begin{scope}[shift={(9.0,0)}]
    \draw[->,hmtgray] (-2.8,0)--(2.8,0) node[right]{\(Q\)};
    \draw[->,hmtgray] (0,-2.0)--(0,2.0) node[above]{\(P\)};
    \fill[hmtlightgold] (0,0)--(2.5,0)
      arc[start angle=0,end angle=57.2958,x radius=2.5cm,y radius=1.35cm]--cycle;
    \draw[hmtviolet,very thick] (0,0) ellipse (2.5cm and 1.35cm);
    \draw[hmtgold,thick] (0,0)--(2.5,0);
    \draw[hmtgold,thick] (0,0)--({2.5*cos(57.2958)},{1.35*sin(57.2958)});
    \node[text=hmtviolet,font=\bfseries] at (0,2.55) {action ellipse};
    \node[align=center] at (.7,.48) {\(\theta=1\)\\area \(\hbar\)};
    \node[align=center,text=hmtgray] at (0,-2.4) {vuelta \(2\pi\): area \(h\)};
  \end{scope}

  \node[box,draw=hmtblue,fill=hmtlightblue,minimum width=3.2cm] (cyl) at (3.2,-4.3)
    {cilindro \(\varepsilon=+1\)\\\(2\pi\mapsto h\)};
  \node[box,draw=hmtviolet,fill=hmtlightviolet,minimum width=4.2cm] (mob) at (9.0,-4.3)
    {Möbius case \(\varepsilon=-1\)\\\(2\pi\mapsto h/2\), \(4\pi\mapsto h\)};
  \draw[arr,hmtgray] (cyl)--node[
    above,align=center,text width=1.9cm,fill=white,inner sep=1pt
  ]{change of\\holonomy} (mob);
  \node[align=center,text=hmtgray,font=\footnotesize] at (6.1,-5.55)
    {The internal phase transports \(\hbar\) per radian;\\the factor two belongs to the visible turn.};
\end{tikzpicture}
}
\caption{The ellipse of rates fixes the form; its lift to the action line fixes
the area. The cylindrical closure and the double Möbius return distinguish the internal phase
from the visible turn.}
\label{fig:elipse-accion-holonomia-rev8}
\end{figure}

\subsection{Hilbert--Beltrami--Klein Geometry on the Action Ellipse}
\label{sec:hilbert-beltrami-klein-elipse-accion}

The same ellipse that carries the symplectic area admits an intrinsic projective
geometry. Let
\[
 \mathbb E_{\rm act}
 :=\left\{(Q,P)\in\mathbb R^2:
 \frac{Q^2}{a_S^2}+\frac{P^2}{b_S^2}<1\right\},
 \qquad
 \mathcal L_{\rm act}(Q,P)
 :=\left(\frac{Q}{a_S},\frac{P}{b_S}\right).
\]
The linear map \(\mathcal L_{\rm act}\) identifies
\(\mathbb E_{\rm act}\) with the unit disk \(\mathbb D\). For two
distinct points \(x,y\in\mathbb E_{\rm act}\), let \(a,b\in\partial
\mathbb E_{\rm act}\) be the intersections of the line \(xy\) with the boundary,
ordered as \(a,x,y,b\). We define
\begin{equation}
 d_{\rm H}^{\rm act}(x,y)
 :=\frac12\log
 \frac{\lVert y-a\rVert\,\lVert x-b\rVert}
      {\lVert x-a\rVert\,\lVert y-b\rVert}.
 \label{eq:distancia-hilbert-elipse-accion}
\end{equation}
The expression is the cross-ratio of the four collinear points and is therefore
independent of the auxiliary norm used on that line.

\begin{proposition}[Beltrami--Klein realization of the action ellipse]
\label{prop:beltrami-klein-elipse-accion}
The pair \((\mathbb E_{\rm act},d_{\rm H}^{\rm act})\) is isometric, through
\(\mathcal L_{\rm act}\), to the Beltrami--Klein model of the hyperbolic plane.
Its unoriented geodesics are exactly the open chords of the
ellipse. The boundary \(\partial\mathbb E_{\rm act}\) simultaneously preserves, in a
typed manner, the symplectic orbit
\[
 \gamma_{\rm act}(\theta)
 =(a_S\cos\theta,b_S\sin\theta),
 \qquad
 \frac12\int_0^\theta(QP'-PQ')\,dt
 =\hbar_{\rm HMT}^{\rm ret}\theta.
\]
\end{proposition}

\begin{proof}
Projective transformations preserve cross-ratios. Since
\(\mathcal L_{\rm act}\) is an invertible linear transformation that carries
the ellipse to the unit disk, it transports
\eqref{eq:distancia-hilbert-elipse-accion} to the Hilbert distance of the
disk. This is the Beltrami--Klein metric, whose geodesics are its chords.
The symplectic identity of the boundary is
\cref{thm:area-fase-elipse-rev8} written in the normalized chart.
\end{proof}

The complete chart is thereby typed by
\[
 \mathfrak R_{\rm act}:
 (A,C^*,\hbar_{\rm HMT}^{\rm ret})
 \longmapsto
 (\mathbb E_{\rm act},d_{\rm H}^{\rm act},
  dP\wedge dQ,\gamma_{\rm act}).
\]
The projective metric governs the interior; the symplectic form and the area
law govern the traversed boundary. A realization as a physical region or
as a domain of covariance matrices is formulated through a subsequent functor
that explicitly preserves these two structures; the preceding construction
already fixes its mathematical domain and the invariants that this functor must preserve.

\endgroup
\begingroup
\let\chapter\section
\let\section\subsection
\let\subsection\subsubsection
\let\subsubsection\paragraph
\let\clearpage\relax
\subsection[Planck's Self-Dual Identity on the Periodicity of Order 108]{Planck self-dual identity in the circular realization of the temporal endomorphism of order 108}
\label{sec:identidad-autodual-planck-ct108-rev3}
\index{Planck!self-dual identity}
\index{calendar of order 108! circular realization}
\index{Compton!reduced and complete readings}

The action line constructed in the preceding section admits a
circular realization compatible with temporal periodicity of order 108.
This realization coordinates three objects that already possess their own types: the
action section, the geometry of one revolution, and the gravitational
transduction. The resulting coordination provides an internal Planck
identity and prepares the positive line on which the general mass
operator will subsequently act.

\subsubsection{Oriented action letters and declared circular realization}

Let
\[
 \mathcal L_S^+,\quad \mathcal L_T^+,\quad \mathcal L_C^+,
 \quad \mathcal L_L^+,\quad \mathcal L_M^+,\quad \mathcal L_G^+
\]
the positive half-lines of action, time, velocity, length, and mass.
The pre-return and post-return coordinates are written
\begin{equation}
 \hbar_{\rm pre}=H_5\,10^{-34}\mathcal U_S,
 \qquad
 \hbar_{\rm ret}=\eta_{\rm ret}\,10^{-34}\mathcal U_S,
 \qquad
 R_{\rm act}=\frac{\hbar_{\rm pre}}{\hbar_{\rm ret}}
 =\frac{H_5}{\eta_{\rm ret}}.
 \label{eq:planck-ct108-cartas-pre-ret-rev3}
\end{equation}
The subscript \(a\in\{\mathrm{pre},\mathrm{ret}\}\) shall designate either
of the two charts. In each one the exact distinction is preserved
\begin{equation}
 h_a=2\pi\hbar_a .
 \label{eq:planck-ct108-h-hbar-rev3}
\end{equation}

Let us fix an elementary duration \(t_0\in\mathcal L_T^+\), an internal
velocity \(c_{\rm int}\in\mathcal L_C^+\), and
\(\ell_0=c_{\rm int}t_0\). The declared circular realization of the
temporal periodicity of order 108 is determined by
\begin{equation}
 T_{108}=108t_0,
 \qquad
 \omega_{108}=\frac{2\pi}{T_{108}},
 \qquad
 R_{108}=\frac{c_{\rm int}}{\omega_{108}}
 =\frac{54}{\pi}\ell_0,
 \label{eq:planck-ct108-realizacion-circular-rev3}
\end{equation}
and by the length of the full turn
\begin{equation}
 C_{108}=2\pi R_{108}=c_{\rm int}T_{108}=108\ell_0.
 \label{eq:planck-ct108-perimetro-rev3}
\end{equation}
Thus, \(R_{108}\) is the radius of the realization and \(C_{108}\) its
perimeter. The angular frequency, the cyclic frequency and the two
action constants are related through
\(\omega_{108}=2\pi/T_{108}\) and \(h_a=2\pi\hbar_a\).

\subsubsection{Correspondence between the order-108 periodicity and Compton readings}

In the chart \(a\), the circular quantum and its mass section are
\begin{equation}
 E_{108,a}=\hbar_a\omega_{108},
 \qquad
 m_{108,a}=\frac{E_{108,a}}{c_{\rm int}^{2}}
 =\frac{\hbar_a\omega_{108}}{c_{\rm int}^{2}}.
 \label{eq:planck-ct108-cuanto-circular-rev3}
\end{equation}

\begin{theorem}[Identity between circular periodicity and Compton readings]
\label{thm:planck-ct108-compton-rev3}
For each chart \(a\in\{\mathrm{pre},\mathrm{ret}\}\), the reduced and complete
Compton readings satisfy
\begin{equation}
 \boxed{
 \bar\lambda_{\rm C}(m_{108,a})
 =\frac{\hbar_a}{m_{108,a}c_{\rm int}}=R_{108},
 \qquad
 \lambda_{\rm C}(m_{108,a})
 =\frac{h_a}{m_{108,a}c_{\rm int}}=C_{108}.}
 \label{eq:planck-ct108-identidad-compton-rev3}
\end{equation}
\end{theorem}

\begin{proof}
Substituting
\eqref{eq:planck-ct108-cuanto-circular-rev3} produces
\[
 \frac{\hbar_a}{m_{108,a}c_{\rm int}}
 =\frac{c_{\rm int}}{\omega_{108}}=R_{108}.
\]
The second equality follows by applying simultaneously
\eqref{eq:planck-ct108-h-hbar-rev3} and
\eqref{eq:planck-ct108-perimetro-rev3}.
\end{proof}

The theorem is homogeneous in
\((\hbar_a,c_{\rm int},t_0)\): the same periodicity determines the radius,
the perimeter, and the two Compton charts. The circular realization is the
declared geometric primitive; the preceding equalities are its
exact consequences.

\subsubsection{Circular selector and internal Planck unit}

The gravitational line is coordinated with the chart \(a\) through the homogeneous
transducer
\begin{equation}
 G_a(\theta)=\theta\,
 \frac{c_{\rm int}^{5}t_0^2}{\hbar_a}
 \in\mathcal L_G^+,
 \qquad \theta\in\mathbb R_{>0}.
 \label{eq:planck-ct108-transductor-rev3}
\end{equation}
Its reduced gravitational radius, evaluated on the circular quantum, is
\begin{equation}
 r_{{\rm g},a}(\theta)
 :=\frac{G_a(\theta)m_{108,a}}{c_{\rm int}^{2}}
 =\theta\frac{\pi}{54}\ell_0.
 \label{eq:planck-ct108-radio-gravitatorio-rev3}
\end{equation}

\begin{proposition}[Planck Circular Selector]
\label{prop:selector-circular-planck-ct108-rev3}
Within the circular realization
\eqref{eq:planck-ct108-realizacion-circular-rev3}, the condition
\(r_{{\rm g},a}(\theta)=R_{108}\) selects a unique positive coordinate,
\begin{equation}
 \boxed{\theta_{108}^{\rm circ}=\left(\frac{54}{\pi}\right)^2.}
 \label{eq:theta-circular-planck-ct108-rev3}
\end{equation}
The coordinate is independent of the chart oriented \(a\).
\end{proposition}

\begin{proof}
The equations
\eqref{eq:planck-ct108-realizacion-circular-rev3} and
\eqref{eq:planck-ct108-radio-gravitatorio-rev3} reduce the condition to
\[
 \theta\frac{\pi}{54}\ell_0=\frac{54}{\pi}\ell_0.
\]
The positivity of \(\ell_0\) and the strict monotonicity of the map
\(\theta\mapsto r_{{\rm g},a}(\theta)\) provide the solution and its
uniqueness.
\end{proof}

Let \(G_{108,a}^{\rm circ}:=G_a(\theta_{108}^{\rm circ})\). The internal units
associated with this chart satisfy
\begin{align}
 \ell_{{\rm P},a}^{\rm circ}
 &:=\sqrt{\frac{\hbar_a G_{108,a}^{\circ}}{c_{\rm int}^{3}}}
 =R_{108},
 &
 t_{{\rm P},a}^{\circ}
 &:=\sqrt{\frac{\hbar_a G_{108,a}^{\circ}}{c_{\rm int}^{5}}}
 =\omega_{108}^{-1},
 \label{eq:planck-ct108-longitud-tiempo-rev3}\\
 m_{{\rm P},a}^{\circ}
 &:=\sqrt{\frac{\hbar_a c_{\rm int}}{G_{108,a}^{\circ}}}
 =m_{108,a},
 &
 E_{{\rm P},a}^{\circ}
 &:=m_{{\rm P},a}^{\circ}c_{\rm int}^{2}=E_{108,a}.
 \label{eq:planck-ct108-masa-energia-rev3}
\end{align}
The equality of radii is therefore published as
\begin{equation}
 R_{108}=\bar\lambda_{\rm C}=r_{\rm g}=\ell_{\rm P}^{\rm circ},
 \qquad
 C_{108}=\lambda_{\rm C}=2\pi r_{\rm g}=2\pi\ell_{\rm P}^{\rm circ}.
 \label{eq:planck-ct108-tres-lectores-rev3}
\end{equation}

The convention of this theorem is the reduced gravitational radius
\(r_{\rm g}=Gm/c_{\rm int}^{2}\). The Schwarzschild radius satisfies
\(r_{\rm s}=2r_{\rm g}\); if this second convention is adopted, its crossing with
\(\bar\lambda_{\rm C}\) occurs at
\(m=m_{{\rm P},a}^{\circ}/\sqrt2\). The declaration of the radius therefore forms
part of the chart and prevents a factor of two from being inadvertently transferred
between the two identities.

The coordinate \(\theta_{108}^{\rm circ}\) belongs to the declared circular
specialization. The general gravitational selector, formulated on the
enriched holonomy and its spectral response, retains its residence in
\cref{sec:selector-gravitatorio-rev11}. Thus, the circular
identity fixes an internal Planck normalization and the general criterion
fixes the conditions that the physical realization of the gravitational
section must satisfy.

\subsubsection{Pre-return/post-return covariance}

The product of the reciprocal sections of action and gravity preserves the
circular geometry:
\begin{equation}
 \frac{m_{108,\rm pre}}{m_{108,\rm ret}}=R_{\rm act},
 \qquad
 \frac{G_{108,\rm pre}^{\circ}}{G_{108,\rm ret}^{\circ}}
 =R_{\rm act}^{-1},
 \qquad
 \hbar_a G_{108,a}^{\circ}
 =\theta_{108}^{\rm circ}c_{\rm int}^{5}t_0^2.
 \label{eq:planck-ct108-covariancia-pre-ret-rev3}
\end{equation}
In particular, \(\ell_{{\rm P},\rm pre}^{\rm circ}\) and
\(\ell_{{\rm P},\rm ret}^{\rm circ}\) coincide with \(R_{108}\). The
bidirectional memory resides in the quotient of the masses and in its
gravitational reciprocal; the geometric scale resides in the invariant product.

Reciprocity admits a symmetric parametrization that separates the common
scale from the orientation of the return. Define
\begin{equation}
 \begin{aligned}
 \hbar_{\rm geom}
 &:=\sqrt{\hbar_{\rm pre}\hbar_{\rm ret}},
 &
 G_{108,\rm geom}^{\circ}
 &:=\sqrt{G_{108,\rm pre}^{\circ}G_{108,\rm ret}^{\circ}},
 &
 \xi_{\rm act}&:=\frac12\log R_{\rm act}.
 \end{aligned}
 \label{eq:planck-ct108-parametrizacion-simetrica-rev3}
\end{equation}
Then
\[
 \hbar_{\rm pre}=\hbar_{\rm geom}e^{\xi_{\rm act}},
 \qquad
 \hbar_{\rm ret}=\hbar_{\rm geom}e^{-\xi_{\rm act}},
\]
whereas the conjugate gravitational sections satisfy
\[
 G_{108,\rm pre}^{\circ}
 =G_{108,\rm geom}^{\circ}e^{-\xi_{\rm act}},
 \qquad
 G_{108,\rm ret}^{\circ}
 =G_{108,\rm geom}^{\circ}e^{\xi_{\rm act}}.
\]
The coordinate \(\xi_{\rm act}\) preserves the pre-return/post-return
orientation; the product \(\hbar_aG_{108,a}^{\circ}\)
preserves geometric scale. This parametrization of the action charts
does not by itself introduce distinct propagation speeds.

\subsubsection{Split self-dual mass line}

Given a chart \(a\), let the split algebra be
\[
 \mathbb D_{\rm sp}:=\mathbb R[j]/(j^2-1),
 \qquad
 P_\pm:=\frac{I\pm j}{2},
\]
and consider its spectral decomposition
\[
 \mathbb D_{\rm sp}
 =\mathbb R P_+\oplus\mathbb R P_-,
 \qquad
 P_\pm^2=P_\pm,\quad
 P_+P_-=0,\quad
 P_++P_-=I.
\]
The broken conjugation swaps the projectors,
\[
 \overline{uP_++vP_-}=vP_++uP_-,
\]
and its norm is
\[
 N_{\rm sp}(uP_++vP_-):=uv.
\]

For \(m\in\mathcal L_M^+\), define the section
\begin{equation}
 Z_a(m)
 :=\frac{m_{{\rm P},a}^{\circ}}{m}P_+
   +\frac{m}{m_{{\rm P},a}^{\circ}}P_- .
 \label{eq:planck-ct108-seccion-partida-masa-rev3}
\end{equation}
The unit \(m_{{\rm P},a}^{\circ}\) also determines the involution
\begin{equation}
 \iota_{{\rm P},a}(m)
 =\frac{(m_{{\rm P},a}^{\circ})^2}{m}.
 \label{eq:planck-ct108-involucion-masas-rev3}
\end{equation}

\begin{theorem}[Split self-dual mass line]
\label{thm:planck-ct108-recta-partida-masa-rev3}
For every \(m\in\mathcal L_M^+\),
\begin{equation}
 \boxed{
 N_{\rm sp}(Z_a(m))=1,
 \qquad
 \overline{Z_a(m)}
 =Z_a\!\left(\iota_{{\rm P},a}(m)\right).}
 \label{eq:planck-ct108-norma-conjugacion-rev3}
\end{equation}
Conjugation has a unique positive fixed point,
\[
 m=m_{{\rm P},a}^{\circ},
 \qquad
 Z_a(m_{{\rm P},a}^{\circ})=I.
\]
\end{theorem}

\begin{proof}
The norm is the product of the two reciprocal coordinates:
\[
 \frac{m_{{\rm P},a}^{\circ}}m
 \frac m{m_{{\rm P},a}^{\circ}}=1.
\]
Conjugation interchanges both coordinates, which is equivalent to replacing
\(m\) by \((m_{{\rm P},a}^{\circ})^2/m\). The positive fixed point satisfies
\(m^2=(m_{{\rm P},a}^{\circ})^2\).
\end{proof}

With the speed
\begin{equation}
 x_a(m)=\log\frac{m}{m_{{\rm P},a}^{\circ}},
 \label{eq:planck-ct108-rapidez-masa-rev3}
\end{equation}
the section and its conjugate are written
\[
 Z_a(m)=e^{-x_a}P_++e^{x_a}P_-,
 \qquad
 \overline{Z_a(m)}
 =e^{x_a}P_++e^{-x_a}P_-.
\]
Their even and odd readers are
\begin{equation}
 Q_a(m):=\frac{e^{x_a}+e^{-x_a}}2=\cosh x_a,
 \qquad
 A_a(m):=\frac{e^{x_a}-e^{-x_a}}2=\sinh x_a.
 \label{eq:planck-ct108-lectores-par-impar-rev3}
\end{equation}
The first preserves the distance to the self-dual point and the second preserves its orientation.

The reciprocal lattice vectors associated with the same coordinate are
\[
 r_{{\rm g},a}(m)=\frac{G_{108,a}^{\circ}m}{c_{\rm int}^{2}},
 \qquad
 \bar\lambda_{{\rm C},a}(m)=\frac{\hbar_a}{mc_{\rm int}},
\]
and satisfy
\begin{equation}
 \frac{r_{{\rm g},a}(m)}{\bar\lambda_{{\rm C},a}(m)}
 =e^{2x_a(m)},
 \qquad
 \frac{\bar\lambda_{{\rm C},a}(m)}{r_{{\rm g},a}(m)}
 =e^{-2x_a(m)}.
 \label{eq:planck-ct108-lectores-exponenciales-rev3}
\end{equation}
Conjugation interchanges the two readers; Planck's identity is their unitary equilibrium.

\begin{exactbox}[title={Planck identity and electron closure}]
The identity \(m=m_{{\rm P},a}^{\circ}\) fixes the self-dual origin of the positive line of masses. The electron section fixes a distinct position on that same line through its own genealogy. The \cref{sec:ley-general-masa-accion-autonoma} constructs the universal operator that transports absolute scale, character, towers, and bidirectional reading; the \cref{sec:cierre-electron-two-way-rev11} then realizes its central rank-one reduction. Thus the Planck unit provides the identity of the coordinate, and the electron provides the first complete material transition on it.
\end{exactbox}

\begin{proposition}[Genealogical realization of the electronic center]
\label{prop:planck-ct108-realizacion-centro-electron-rev3}
Let
\[
 F_e^+=\mathbb R_{>0}(P_++P_-)
\]
the positive half-line of the rank-one central fiber, let
\(\mathcal G_e\) be the space of enriched electronic
genealogies and let \(\mathcal L_M^+\) be the positive mass
line. A realization morphism
\[
 \mathfrak R_e:
 F_e^+\times\mathcal G_e\longrightarrow\mathcal L_M^+
\]
can send the identity \(I=P_++P_-\) to a mass different from
\(m_{{\rm P},a}^{\circ}\) without altering the
\cref{thm:planck-ct108-recta-partida-masa-rev3}: the Planck identity fixes the
autodual origin of the chart, while the electronic genealogy fixes a
displacement on that chart.
\end{proposition}

\begin{proof}
If \(\chi_e:\mathcal G_e\to\mathbb R_{>0}\) is the positive character constructed
by the electronic record, the rule
\[
 \mathfrak R_e(\lambda I,g)
 =\lambda\,m_{{\rm P},a}^{\circ}\chi_e(g)
\]
is positive and multiplicative in the genealogical coordinate. For the neutral
genealogy, \(\chi_e=1\), Planck's origin is recovered; for a genealogy with
\(\chi_e(g)\ne1\), the same internal identity occupies a different position on
the line. The specific electronic closure preserves its source in the
\cref{sec:cierre-electron-two-way-rev11}.
\end{proof}

\paragraph{Domain of the Planck self-dual identity: temporal periodicity of order \(108\), declared circular realization, and separation of the electronic closure}
The pre-return/post-return action line and the temporal periodicity of
order (108) precede the circular realization, which is taken as a declared
geometric primitive. The identity between circular
periodicity and the Compton readings, the
selector \(\theta_{108}^{\rm circ}\), the internal Planck unit, and the
involution of the positive line are exact results once that
realization has been fixed. General gravitational selection and electronic
closure are constructed separately in the indicated sections.

\endgroup
\begingroup
\let\chapter\section
\let\section\subsection
\let\subsection\subsubsection
\let\subsubsection\paragraph
\let\clearpage\relax
\section{The general law before its species-specific realizations}
\label{sec:ley-general-masa-accion-autonoma}
\index{general operatorial mass law}
\index{action!absolute mass scale}
\index{spectral towers!global mass law}

Dimensional Mechanics receives from HMT two branches already constructed. The first
provides an action line with two oriented coordinates; the second
provides the genealogy of each persistent section. Mass appears upon
making the genealogical character act on a section of the mass line.
This confluence precedes every particular case and, accordingly, the
scale \(10^{-34}\) belongs to the general law rather than to an isolated species.

\subsection{Absolute Branch: From the Local Determinant to the Mass Line}

The determinantal reader fixes locally
\begin{equation}
 \Sigma_H=\varphi_{\rm HMT}\alpha_{\rm HMT}^{16},
 \qquad \varepsilon_H:=\left\lfloor\log_{10}\Sigma_H\right\rfloor=-34 .
 \label{eq:ley-masa-sigma-menos34-autonoma}
\end{equation}
The power sixteen follows from the cokernel
\(
 \mathbb Z/2\oplus\mathbb Z/2\oplus\mathbb Z/4
\)
of the determinantal block. On the internal action line
\(\mathcal L_S\) the two sections are obtained
\begin{equation}
 \boldsymbol\hbar_{\rm pre}
 =H_5\,10^{-34}\mathcal U_S,
 \qquad
 \boldsymbol\hbar_{\rm ret}
 =\eta_{\rm ret}\,10^{-34}\mathcal U_S,
 \label{eq:ley-masa-hbar-pre-ret-autonoma}
\end{equation}
and the additive and multiplicative letters of its unfolding are
\begin{equation}
 \delta_{2w}=H_5-\eta_{\rm ret},
 \qquad
 R_{\rm act}
 =\frac{\boldsymbol\hbar_{\rm pre}}
        {\boldsymbol\hbar_{\rm ret}}
 =\frac{H_5}{\eta_{\rm ret}}>0 .
 \label{eq:ley-masa-razon-accion-autonoma}
\end{equation}

Let \(t_0\in\mathcal L_T^+\) be the elementary duration of the TPK calendar and
\(c_{\rm int}\in\mathcal L_C^+\) its internal speed. The periodicity of
108 iterations of the temporal endomorphism CT108 fixes the angular frequency
\(\omega_0=2\pi/(108t_0)\). The absolute sections preceding and following the
return are, therefore,
\begin{equation}
 \begin{aligned}
 \mathbf m_\star^{\rm pre}
 &=\boldsymbol\hbar_{\rm pre}\,\omega_0c_{\rm int}^{-2},\\
 \mathbf m_\star^{\rm ret}
 &=\boldsymbol\hbar_{\rm ret}\,\omega_0c_{\rm int}^{-2},
 \end{aligned}
 \qquad
 \frac{\mathbf m_\star^{\rm pre}}
      {\mathbf m_\star^{\rm ret}}=R_{\rm act}.
 \label{eq:secciones-absolutas-masa-autonoma}
\end{equation}
Both belong to
\(
 \mathcal L_M=\mathcal L_S\otimes\mathcal L_T^{-1}
 \otimes\mathcal L_C^{-2}
\).
Thus the decade, which fixes the absolute scale, and the
quotient of the pre-return/post-return bidirectional splitting, which
transports the memory of the closed evolution, are precisely separated.

In positive internal coordinates
\(t_0=\widehat t_0\mathcal U_T\) and
\(c_{\rm int}=\widehat c\mathcal U_C\), the scalar reduction of any
section whose physical realization has already been fixed takes the form
\begin{equation}
 \mathbf m_X^\beta
 =10^{-34}\eta_{\rm ret}\,
  \frac{2\pi}{108\widehat t_0}\,\widehat c^{-2}
  a_X^{(0)}R_{\rm act}^{\kappa_X}\,\mathcal U_M .
 \label{eq:ley-escalar-general-masas-autonoma}
\end{equation}
The coordinate \(a_X^{(0)}\) proceeds of the character and of the hierarchy harmonic;
\(\kappa_X\) proceeds of the reader of route. In particular, the decade
\(-34\) belongs to the section absolute and does not  incorporates to
\(\kappa_X\).

\subsection{Genealogical branch: record, character, and harmonic hierarchy}

Each surviving section traverses the chain
\begin{equation}
 \Omega_{\rm HMT}^{\rm surv}
 \longrightarrow\Gamma_{108}^{\rm enr}
 \longrightarrow\mathcal L_{108}
 \longrightarrow\sigma\in\mathbb Z^5
 \longrightarrow(\sigma,\eta)
 \longrightarrow\chi_{\rm full}(\sigma,\eta)\in\mathbb R_{>0}.
 \label{eq:cadena-genealogica-ley-masa-autonoma}
\end{equation}
The record preserves the order of events; the signature preserves their five
integer invariants; the character evaluates them without erasing their provenance. Its
spectral prolongation lies in the full semigroup
\begin{equation}
 \mathfrak S=\langle90,120\rangle
 =\{90,120\}\cup\{30r:r\ge6\},
 \qquad
 s_h=q_-^h-q_+^h,
 \quad c_h=q_-^h+q_+^h,
 \label{eq:semigrupo-global-masas-autonoma}
\end{equation}
and adopts the form
\begin{equation}
 \chi_{\rm full}(\sigma,\eta)
 =\chi_0(\sigma)
  \exp\!\left(\sum_{h\in\mathfrak S}\eta_hs_h\right).
 \label{eq:caracter-completo-masas-autonoma}
\end{equation}
The heights \(90,120,180,210,240,270,360,420\) are eight especially visible
genealogical witnesses. The law belongs to all of
\(\mathfrak S\); the second-order recurrence of \((s_h,c_h)\) extends
the witnesses to the complete harmonic hierarchy.

\subsection{Universal operator on the \(13\) multisections of the material domain}

The atlas provides \(81\) positions, \(56\) route families, \(13\)
multisections, and \(324\) oriented routes. For a multisection \(F\), let
\(\mathcal H_F=\bigoplus_{\Gamma\in F}\mathbb C e_\Gamma\). The basal character
and the towers define the positive operator
\begin{equation}
 \widehat{\mathcal A}^{(0)}_F e_\Gamma
 =\chi_{\rm full}(\sigma_\Gamma,\eta_\Gamma)e_\Gamma.
 \label{eq:operador-caracter-basal-masas-autonoma}
\end{equation}
The two histories preserved by TPK induce readers
\(r\in\{D,\Omega\}\): \(K_D\) acts on the chronological record of
108 iterations of the CT108 temporal endomorphism, and \(K_\Omega\) on the signed
dodecaphase word. If
\(\widehat K_F^r e_\Gamma=\kappa_r(\Gamma)e_\Gamma\), let
\(F_s\subseteq F\) be the set of routes of the section or physical block \(s\), and
define
\begin{equation}
 \mathcal H_{F,s}:=\bigoplus_{\Gamma\in F_s}\mathbb C e_\Gamma,
 \qquad
 \widehat{\mathcal A}^{(0)}_{F,s}
 :=\left.\widehat{\mathcal A}^{(0)}_F\right|_{\mathcal H_{F,s}},
 \qquad
 \widehat B_{F,s}^r
 :=\left.\widehat K_F^r\right|_{\mathcal H_{F,s}}.
 \label{eq:restricciones-bloque-ley-masas-autonoma}
\end{equation}
The basis diagonal of routes used by the certificates current makes invariant
\(\mathcal H_{F,s}\) under both operators and converts
\(\widehat B_{F,s}^r\) in a restriction self-adjoint. With these types, the law
general is
\begin{equation}
 \boxed{
 \widehat{\mathbf M}^{\beta,r}_{F,s}
 =\mathbf m_\star^{\rm ret}\otimes
  R_{\rm act}^{\widehat B_{F,s}^r/2}
  \widehat{\mathcal A}^{(0)}_{F,s}
  R_{\rm act}^{\widehat B_{F,s}^r/2},
 \qquad r\in\{D,\Omega\}.}
 \label{eq:ley-operatorial-general-masas-autonoma}
\end{equation}
The symmetric writing preserves positivity without presupposing commutation. In the
simultaneously diagonal basis employed by the
current certificates, it reduces to
\begin{equation}
 \widehat{\mathbf M}^{\beta,r}_F
 =\mathbf m_\star^{\rm ret}\otimes
  \widehat{\mathcal A}^{(0)}_F
  \exp\!\left((\log R_{\rm act})\widehat K_F^r\right).
 \label{eq:ley-operatorial-diagonal-masas-autonoma}
\end{equation}
The equation preserves, in a single object, the absolute scale, the global
character, the entire hierarchy of towers, and the route-dependent bidirectional
corrector. The physical chart of each sector subsequently acts on the spectrum of
\eqref{eq:ley-operatorial-general-masas-autonoma}. The null sector has its
own morphism to zero; the composite registers are concatenated before
the signature is evaluated; the topological sectors first publish their fusion
and braiding invariants.

\begin{theorem}[Covariance and Cancellation of the Common Scale]
\label{thm:covariancia-ley-general-masas-autonoma}
The law~\eqref{eq:ley-operatorial-general-masas-autonoma} is positive and
covariant under unitary change of basis in each fiber. For two self-sections
\(X,Y\) realized in the same mass chart,
\begin{equation}
 \frac{m_Y}{m_X}
 =\frac{\chi_{\rm full}(\sigma_Y,\eta_Y)}
        {\chi_{\rm full}(\sigma_X,\eta_X)}
  R_{\rm act}^{\kappa_r(\Gamma_Y)-\kappa_r(\Gamma_X)}.
 \label{eq:cociente-ley-general-masas-autonoma}
\end{equation}
The common factors \(10^{-34}\mathcal U_S\), \(2\pi/(108t_0)\), and
\(c_{\rm int}^{-2}\) fix the absolute section and cancel exactly in
the quotient.
\end{theorem}

\begin{proof}
The operator \(\widehat{\mathcal A}^{(0)}_F\) is positive because its
eigenvalues belong to \(\mathbb R_{>0}\); the exponential of a
self-adjoint operator is positive. Conjugation by a unitary matrix preserves the
spectrum and therefore the law. Upon dividing two eigenvalues realized in the
same chart, the common section \(\mathbf m_\star^{\rm ret}\) cancels, and the result
is ~\eqref{eq:cociente-ley-general-masas-autonoma}.
\end{proof}

The central electronic fiber has rank one and constitutes the first
scalar corollary of this law. The remaining multisections preserve the pair
of operators, their spectra, and the physical morphism that selects flavor,
color, orientation, or composition. This hierarchy allows the algebraic
result, its dimensional realization, and the external contrast to be published together without
confusing their domains.

\endgroup
\begingroup
\let\chapter\section
\let\section\subsection
\let\subsection\subsubsection
\let\subsubsection\paragraph
\let\clearpage\relax
% Vista científica exacta materializada desde una rama condicional.
% Fuente: source/demostraciones_primarias/09_06h_atlas_familias_rev6.tex; rama: HMTIncludeAtlasStructural.
% No modifica el contenido matemático de la rama de origen.
\chapter[Nonadic classification of families]{Nonadic classification of trajectory families}
\label{chap:atlas-genealogia-masas}
\index{nonadic atlas}
\index{family of routes}
\index{multisection}
\index{null sector}

The atlas does not begin with a list of particles or a table of masses.
It begins with a finite oriented support and with the readers that preserve the
history of its routes.  Its function is to answer a question prior to
any metrological comparison: what classes of persistence can the
nonadic geometry support, and what information must be preserved to distinguish them?

The answer has three cardinals, but not three rival inventories:
\[
 81\ \text{cells},\qquad
 56\ \text{route families},\qquad
 13\ \text{family--level multisections}.
\]
The first number counts oriented positions; the second identifies the
positions that publish the same finite route signature; the third gathers the
positions that share a family and depth relative to the center. None
of them counts known masses, names of species, or rows of an
external comparison.

\section{The complete geometry of the \(81\) cells}

Let
\[
 I_9=\{1,\ldots,9\},
 \qquad
 \mathcal C_{81}=I_9\times I_9.
\]
The coordinates \((r,c)\) preserve the two directions of APP\@. The parity
of each coordinate defines four internal sectors:
\[
 \operatorname{fam}(r,c)=
 \begin{cases}
  \ell^- ,& r\equiv1,\ c\equiv1\pmod2,\\
  \nu ,& r\equiv1,\ c\equiv0\pmod2,\\
  d ,& r\equiv0,\ c\equiv1\pmod2,\\
  u ,& r\equiv0,\ c\equiv0\pmod2.
 \end{cases}
 \tag{A1}\label{eq:familias-paridad-atlas}
\]
Depth in the atlas is not added as an independent label.  It is the
Chebyshev distance to the center:
\[
 G(r,c)=\max\{|r-5|,|c-5|\}.
 \tag{A2}\label{eq:generacion-chebyshev-atlas}
\]
Therefore, family and level follow from the position itself.  In particular,
the center is not obtained by selecting a species from outside: it is the unique
cell of depth zero.

\begin{theorem}[Census of families and generations of the atlas]
\label{thm:censo-atlas-81}
The maps \eqref{eq:familias-paridad-atlas} and
\eqref{eq:generacion-chebyshev-atlas} decompose the atlas in a
disjoint manner as
\[
 81=45_{\rm leptones}+36_{\rm quarks}
   =25_{\ell^-}+20_{\nu}+20_{d}+16_{u}.
 \tag{A3}\label{eq:censo-81-cuatro-familias}
\]
Refinement by \(G\) yields exactly the thirteen multisections
\[
\begin{array}{c|c|c}
\text{family}&\text{levels}&\text{cardinals}\\ \hline
\ell^-&G0,G2,G4&1,8,16\\
\nu&G1,G2,G3,G4&2,4,6,8\\
d&G1,G2,G3,G4&2,4,6,8\\
u&G1,G3&4,12.
\end{array}
\tag{A4}\label{eq:trece-multisecciones-atlas}
\]
The cell \((5,5)\) is the sole component of
\(\mathfrak S_{\ell^-,0}\).
\end{theorem}

\begin{proof}
In \(I_9\) there are five odd coordinates and four even ones.  The four parity
classes therefore have cardinalities
\[
 5\cdot5=25,\qquad
 5\cdot4=20,\qquad
 4\cdot5=20,\qquad
 4\cdot4=16,
\]
which proves \eqref{eq:censo-81-cuatro-familias}; the first two
classes comprise \(45\) cells and the last two, \(36\).

For the refinement, we enlarge squares centered at \((5,5)\).
In the odd--odd class there is one cell at radius zero, \(3^2-1=8\)
in the shell of radius two, and \(5^2-3^2=16\) in the shell of radius four.
In the odd--even class, the rectangles accumulated up to radii
\(1,2,3,4\) contain, respectively,
\[
 1\cdot2=2,\qquad 3\cdot2=6,\qquad
 3\cdot4=12,\qquad5\cdot4=20
\]
cells; their successive differences are \(2,4,6,8\). The
even--odd class gives the same calculation with the coordinates exchanged. In the
even--even class, there are \(2^2=4\) cells at radius one and
\(4^2-2^2=12\) at radius three. These shells are disjoint and exhaust
\(I_9^2\), so no other multisection remains.
\end{proof}

The word \emph{level} denotes in this chapter the combinatorial depth
\(G\).  The subsequent physical reading preserves that geometry but does not
define it.  This prevents particle names from being turned into premises
of the census that must explain them.

\section{Central inversion, \(41\) orbits, and uniqueness of the centre}

The inverse of the atlas is
\[
 \rho_{\rm c}(r,c)=(10-r,10-c).
 \tag{A5}\label{eq:inversion-central-atlas-rev6}
\]
Since \(10-r\) has the same parity as \(r\), inversion preserves the
four families.  Moreover,
\[
 |(10-r)-5|=|r-5|,
 \qquad
 |(10-c)-5|=|c-5|,
\]
and therefore preserves the level \(G\).  Each of the thirteen
multisections is \(\rho_{\rm c}\)-stable.

\begin{proposition}[Orbits and Central Exception]
\label{prop:orbitas-centro-atlas-rev6}
The action of \(\rho_{\rm c}\) on \(\mathcal C_{81}\) has exactly
\(41\) orbits.  Its unique fixed point is \((5,5)\).  Consequently, the
multisection \(\mathfrak S_{\ell^-,0}\) admits the unique pointwise
\(\rho_{\rm c}\)-equivariant selector; in the other twelve, the natural object is the
orbit or the complete multisection.
\end{proposition}

\begin{proof}
The equation
\[
 (r,c)=\rho_{\rm c}(r,c)
\]
is equivalent to \(2r=2c=10\), hence \(r=c=5\).  The other eighty cells are
grouped into forty pairs, whence \(1+40=41\) orbits.

Let \(y\) be a noncentral multisection and suppose that an equivariant point selector
chooses \(a(y)\in\mathfrak S_y\).  Since \(\rho_{\rm c}\) acts
trivially on the family--level type, equivariance would impose
\[
 a(y)=a(\rho_{\rm c}y)=\rho_{\rm c}a(y).
\]
The chosen point would have to be fixed, but the only fixed point lies in
\(\mathfrak S_{\ell^-,0}\).  The contradiction proves the assertion.
\end{proof}

This proposition explains by which a species internal is not, in general, a
cell isolated.  The electron central is exceptional because its multisection
is reduceds to a point.  In the remaining families, the orientation conjugate form
part of the identity that must is preserved.

\section{The \(56\) route families}

Each cell bears the finite shadow of its transport history:
\[
 \mathcal R(r,c)=
 \bigl(
 n_A,\ s_C,\ k_{\Delta_4},\
 \nu_{120},\ \nu_{270},\ \tau_{12}
 \bigr)_{r,c}.
 \tag{A6}\label{eq:lector-familia-ruta-atlas}
\]
The first five coordinates are integers; \(\tau_{12}\) is the oriented
dodecaphase word. Two cells belong to the same route family
if and only if they have the same sextuple \(\mathcal R\).

\begin{theorem}[Census of route families and minimal memory]
\label{thm:censo-56-familias-ruta-rev6}
The image of \(\mathcal R\) contains exactly \(56\) elements.  The
histogram of the sizes of its fibres is
\[
 41\times1,\qquad
 10\times2,\qquad
 2\times3,\qquad
 1\times4,\qquad
 2\times5.
 \tag{A7}\label{eq:histograma-fibras-ruta}
\]
In particular,
\[
 41+10+2+1+2=56,
\]
\[
 41\cdot1+10\cdot2+2\cdot3+1\cdot4+2\cdot5=81.
 \tag{A8}\label{eq:doble-censo-rutas}
\]
Fixing the lexicographic order within each fiber, the rank
\[
 \kappa(r,c)\in\{0,1,2,3,4\}
\]
makes the application injective
\[
 (r,c)\longmapsto\bigl(\mathcal R(r,c),\kappa(r,c)\bigr).
 \tag{A9}\label{eq:lift-memoria-cinco-atlas}
\]
Five is the minimum size of a uniform alphabet capable of performing this lift.
\end{theorem}

\begin{proof}
The exhaustive enumeration of \(I_9^2\) groups the eighty-one sextuples
\eqref{eq:lector-familia-ruta-atlas} by coordinatewise equality.
It produces the five fiber sizes of
\eqref{eq:histograma-fibras-ruta}; the two identities
\eqref{eq:doble-censo-rutas} simultaneously verify the number of classes
and the coverage of the eighty-one cells.

The lexicographic rank distinguishes the elements of each fiber, so that
\eqref{eq:lift-memoria-cinco-atlas} is injective.  Since two fibers
of cardinal five exist, any common alphabet with fewer than five symbols
would identify two cells of one of them by the pigeonhole principle.
\end{proof}

The rank \(\kappa\) is positional memory within a fiber; it is not
charge, mass, or color. Its function is to recover the cell that the visible
route signature had identified with others. The number \(56\) counts classes of the
sextuple \(\mathcal R\), whereas the number \(13\) counts classes of the pair
\((\operatorname{fam},G)\). They are different quotients of the same atlas.

\paragraph{Cellular and orbital quotients of the atlas: domains and projection maps.}
The raw CT--108 projection and the scientific classification read different
data from the same eighty-one cells. If
\(q_{\rm CT}\) preserves only the form of the dodecaphase tritic word,
while \(\mathcal R\) preserves the route sextuple and
\(q=(\operatorname{fam},G)\) preserves family and depth, the typed
table is
\[
\begin{array}{rcll}
 \mathcal C_{81}
 &\xrightarrow{\ q_{\rm CT}\ }&
 \mathcal W^{\rm raw}_{7}
 & \text{raw CT--108 forms},\\[2mm]
 \mathcal C_{81}
 &\xrightarrow{\ \mathcal R\ }&
 \Sigma^{\rm ruta}_{56}
 & \text{route families},\\[2mm]
 \mathcal C_{81}
 &\xrightarrow{\ q\ }&
 \Sigma_{13}
 & \text{family--level multisections}.
\end{array}
\tag{A9a}\label{eq:dos-cocientes-atlas-rev8}
\]
The seven raw fibres have multiplicities
\[
 54,2,2,2,7,7,7,
 \qquad54+3\cdot2+3\cdot7=81.
\]
Therefore, the archaeological census \(81\to7\) does not replace the atlas \(81/56/13\). Nor does the abbreviated notation \(81/56/13\) assert the existence of a canonical arrow \(56\to13\): some fibers of \(\mathcal R\) traverse more than one family--level pair. The maps \(\mathcal R\) and \(q\) are parallel quotients and preserve different invariants. In the quark sector, the active convention remains
\[
 \operatorname{col}(r,c)=r-c\pmod3;
\]
the historical notation \(c-r\) is its conjugate convention, not another color or
another quotient of the atlas.

\section{Distinction between cell, route family, and multisection}

Define
\[
 q:\mathcal C_{81}\longrightarrow\Sigma_{13},
 \qquad
 q(r,c)=\bigl(\operatorname{fam}(r,c),G(r,c)\bigr),
 \tag{A10}\label{eq:cociente-trece-multisecciones}
\]
and, for \(y\in\Sigma_{13}\),
\[
 \mathfrak S_y^{\rm atlas}
 =q^{-1}(y)
 =\{(r,c)\in\mathcal C_{81}:q(r,c)=y\}.
 \tag{A11}\label{eq:multiseccion-atlas-rev6}
\]
The four levels of reading are then separated:

\begin{enumerate}[label=(\roman*)]
 \item a \emph{cell} is an oriented position \((r,c)\), with all its
 local data;
 \item a \emph{route family} is a fibre of \(\mathcal R\): it preserves
 the same finite transport signature;
 \item a \emph{multisection} is a fiber of \(q\): it preserves the family,
 depth, and all its conjugate orientations;
 \item a \emph{physical representation} is the subsequent realization of an
enriched multisection in a state space; it is not another partition of the
board or the arbitrary choice of one of its cells.
\end{enumerate}

This distinction gives mathematical content to the word \emph{species}. The internal type is the multisection. A concrete particle is a persistent section realized within it. The physical representation preserves the action, orientation, and observables of the type; it does not work backward to redefine the atlas from an already known mass.

\section{Color residual as ternary orientation}

In the quark sector, formed by the \(36\) cells with even \(r\), we define
\[
 \operatorname{col}(r,c)=r-c\pmod3,
 \tag{A12}\label{eq:color-residual-atlas-rev6}
\]
with the dictionary
\[
 0\longmapsto R,\qquad1\longmapsto G,\qquad2\longmapsto B.
\]

\begin{proposition}[Residual color balance]
\label{prop:color-atlas}
The thirty-six quark cells are distributed exactly as
\[
 36=12_R+12_G+12_B.
 \tag{A13}\label{eq:balance-color-atlas-rev6}
\]
The central inversion fixes \(R\) and interchanges \(G\) and \(B\).
\end{proposition}

\begin{proof}
Upon fixing any of the four even values of \(r\), as
\(c=1,\ldots,9\) is traversed, each class modulo three appears three times. Each residue
therefore has \(4\cdot3=12\) representatives. Moreover,
\[
 \operatorname{col}\bigl(\rho_{\rm c}(r,c)\bigr)
 =(10-r)-(10-c)
 =-(r-c)\pmod3.
\]
The modular negation fixes the zero class and swaps the one and two classes.
\end{proof}

The tritic \(R,G,B\) is therefore an exact observable of the atlas and not a
decoration added to a table of quarks.  The color representation is
subsequently constructed on this ternary support; the census \(12+12+12\) is its
prior geometry.

\section{From the surviving extension to the signature}

The atlas classifies finite cuts.  The identity of a particle requires the
deep history from which that cut proceeds.  Let
\[
 \Omega_{\rm HMT}^{\rm surv}
 =\varprojlim_m\mathcal S_m
\]
the inverse system of surviving states.  The projection onto the atlas \(\pi_{81}\) and the route
shadow form the chain
\[
 \Omega_{\rm HMT}^{\rm surv}
 \xrightarrow{\ \operatorname{sh}_{108}\ }
 \Gamma_{108}^{\rm enr}
 \xrightarrow{\ \mathcal E\ }
 \mathcal L_{108}
 \xrightarrow{\ L_{\rm tick}\ }
 \mathbb Z^5.
 \tag{A14}\label{eq:cadena-extension-firma-atlas-rev6}
\]
Here \(\operatorname{sh}_{108}\) publishes a route of \(108\) steps without
forgetting sheet, orientation, boundary, or predecessors; \(\mathcal E\) forms the
dodecaphase record; and \(L_{\rm tick}\) extracts
\[
 v(x)=
 \bigl(n_A,s_C,k_{\Delta_4},\nu_{120},\nu_{270}\bigr).
 \tag{A15}\label{eq:firma-cinco-extension-rev6}
\]
None of these three arrows refers to a mass, a unit, or the name of
a particle.

The reason for preserving the ledger before projection is apparent in its twelve
oriented events \(e_0,\ldots,e_{11}\). For \(0\leq i\leq5\), let
\[
 p_i=e_i+e_{i+6},
 \qquad
 a_i=e_i-e_{i+6},
 \qquad
 s_C=\sum_{i=0}^{5}p_i,
\]
and, for \(0\leq i<5\),
\[
 M_i=p_i-p_5.
\]
Then
\[
 p_5=\frac{s_C-\sum_{i=0}^{4}M_i}{6},
 \qquad
 p_i=p_5+M_i,
\]
\[
 e_i=\frac{p_i+a_i}{2},
 \qquad
 e_{i+6}=\frac{p_i-a_i}{2}.
 \tag{A16}\label{eq:reconstruccion-registro-doce}
\]
The identity
\[
 1+5+6=12
\]
is thus a reversible reconstruction: one total, five differences, and six
orientations recover the twelve events.  Projecting onto the five-integer
signature before composing histories would eliminate precisely the placement and
orientation that make it possible to decide whether two routes can be composed.

The quotient \(q\) is lifted to the deep history through projection to the atlas
\(\pi_{81}:\Omega_{\rm HMT}^{\rm surv}\to\mathcal C_{81}\):
\[
 \widetilde{\mathfrak S}_y
 =\{x\in\Omega_{\rm HMT}^{\rm surv}:q(\pi_{81}x)=y\}.
 \tag{A17}\label{eq:multiseccion-enriquecida-rev6}
\]
This is the enriched multisection.  Its shadow belongs to the atlas; its route
belongs to \(\Gamma_{108}^{\rm enr}\); its complete memory belongs to the
record; and its evaluation is applied only at the end.  The genealogy is not a
note appended to the value: it is the object that makes the value possible and preserves the
other readings of the same section.

\section{Null sector bifurcation prior to the positive character}

The mass character acts on the entire signature and has positive codomain:
\[
 \chi_{\rm mass}:\mathbb Z^5\longrightarrow\mathbb R_{>0}.
 \tag{A18}\label{eq:caracter-positivo-atlas-rev6}
\]
By definition, there is no signature \(v\) with
\(\chi_{\rm mass}(v)=0\). The null sector is, consequently, not a
fourteenth multisection nor the degenerate limit of the thirteen preceding ones.

Its algebraic support belongs to the split branch.
\[
 \mathcal A_{\rm sp}
 =\mathbb R[\mathsf R]/(\mathsf R^2-1),
 \qquad
 P_\pm=\frac{1\pm\mathsf R}{2}.
\]
Every element is written in a unique way as
\[
 z=r_+P_++r_-P_-,
 \qquad
 N_{\rm sp}(z)=r_+r_-,
\]
and the two lines
\[
 \mathbb R P_+,\qquad\mathbb R P_-
\]
are null because \(N_{\rm sp}(\lambda P_\pm)=0\). A null section is an
open route transported along one of those directions. The material interior
requires the coupling of both, for which \(r_+r_->0\).

Thus are typed two branches:
\[
 \text{material persistence}
 \longrightarrow\mathbb Z^5
 \xrightarrow{\chi_{\rm mass}}\mathbb R_{>0},
\]
\[
 \text{null propagation}
 \longrightarrow\partial_{\rm null}\mathcal C_{\rm sp}^{+}.
\]
The photonic and gauge realizations subsequently read the second branch.  Their
separation from the positive character prevents the photon and gluon from being erased because they do not appear
in a list of masses, while also preventing a null mass from being fabricated by means of a
character whose codomain excludes it.

\section{Finite classification system and path memory}

The three quotients of the chapter fulfill complementary functions:
\[
 \mathcal C_{81}
 \xrightarrow{\ \mathcal R\ }\mathcal R_{56},
 \qquad
 \mathcal C_{81}
 \xrightarrow{\ q\ }\Sigma_{13},
 \qquad
 \mathcal C_{81}
 \xrightarrow{\ /\langle\rho_{\rm c}\rangle\ }
 \mathcal O_{41}.
\]
The first preserves the route signature; the second preserves family and level;
the third preserves central conjugation.  Together with residual color and the
lift \(\kappa\), they form a finite system capable of reconstructing
which information is seen, which is identified, and which must remain as
memory.

The causal chain of the atlas is thereby fixed:
\[
 \boxed{
 \begin{gathered}
 \text{surviving extension}
 \longrightarrow\text{enriched route}
 \longrightarrow\text{record}
 \longrightarrow\text{signature}\\
 \longrightarrow\text{character}
 \longrightarrow\text{dimensional reading}
 \end{gathered}}
 \tag{A19}\label{eq:cadena-rectora-atlas-rev6}
\]
The eighty-one cells constitute the local domain; the fifty-six
families express route coincidences; the thirteen multisections are the
internal types; the forty-one orbits preserve conjugation; and the
center \((5,5)\) is the sole pointwise exception. Mass appears thereafter
as a reading of an already constructed persistence. The atlas is not an
abridged table of results: it is the structure explaining why a table can
exist.

\endgroup
\begingroup
\let\chapter\section
\let\section\subsection
\let\subsection\subsubsection
\let\subsubsection\paragraph
\let\clearpage\relax
\section{From APP seeds to mixture and mass registration}
\label{sec:semillas-registro-mezcla-autonoma}
\index{seed!factorization sum--product}
\index{mixture!origin in the seeds}

The physical law receives a domain that has already preserved the incompatibility of
the two APP evaluations. If
\(\mathcal S_+\) and \(\mathcal S_\times\) denote the oriented seeds of
the additive and multiplicative leaves, then
\begin{equation}
 \Omega_{\rm seed}=\mathcal S_+\times\mathcal S_\times,
 \qquad
 |\Omega_{\rm seed}|=324^2=104,976.
 \label{eq:dominio-semillas-md-autonoma}
\end{equation}
The exact factorization of the emitter produces
\begin{equation}
 104,976
 \longrightarrow18\times26=468\ U_6
 \longrightarrow243\ w_6\subset\mathbb F_3^6,
 \qquad |\mathbb F_3^6|=729.
 \label{eq:embudo-semillas-md-autonoma}
\end{equation}
The four cardinalities count distinct objects: initial conditions of
two sheets, decimal emissions, published ternary words, and the algebraic
ambient space. In addition to its emission, each seed preserves multiplicity,
diagonal, sheet, orbit, emission signature, residues modulo three and modulo
nine, turns, sheet changes, endpoint, displacements, and cardinal
counts. Those data constitute the vocabulary of the subsequent register.

The factorization organizes six microfibers of cardinality \(1008\), six phase
Fano planes, thirty-six fibers of cardinality \(28\), and the band
\(1944=27\cdot72\), diagonalized by \((\mathbb Z/3)^3\). Upon averaging the
nonadic word of the TPK, the conditional projectors of the two sheets
define the canonical kernel
\begin{equation}
 K_{\rm cat}=\frac13(I+E_++E_\times),
 \qquad
 \operatorname{Spec}(K_{\rm cat})
 =\left\{1^{[1]},
 \left(\frac23\right)^{[42]},
 \left(\frac13\right)^{[425]}\right\}.
 \label{eq:kernel-canonico-semillas-md-autonoma}
\end{equation}
The TPK phase selects the active renewal and preserves the nine
spatial bands; the average~\eqref{eq:kernel-canonico-semillas-md-autonoma}
publishes the stochastic shadow of the complete calendar.

\subsection{Bigrading of the Seeds and \(6\) Incidence Planes}

In each of the six microfibers of cardinality \(1008\), the bipartite graph
retains \(54\) additive vertices and \(56\) multiplicative vertices. The
former decompose into \(36\) vertices of degree \(24\) and \(18\) of
degree \(8\); the latter have degree \(18\). The incidence count
agrees on both sheets:
\begin{equation}
 1008=36\cdot24+18\cdot8=56\cdot18.
 \label{eq:bidegraduacion-microfibra-masas-autonoma}
\end{equation}
This equality fixes the emitter bigrading before projecting it onto
the route families.

The six words
\begin{equation}
 001112,\quad010211,\quad100121,\quad112001,\quad121100,\quad211010
 \label{eq:seis-palabras-microfibra-masas-autonoma}
\end{equation}
form an orbit of \(D_3\simeq S_3\) and have the profile
\(432+4\cdot144=1008\). The quotient \(1008=36\cdot28\) organizes seven
blocks
\[
 Q=\{B_3,B_6,B_9,P_1,P_2,P_3,f\},
\]
where \((B_3,B_6,B_9)\) are the transverse tetrads and
\((P_1,P_2,P_3,f)\), the four lateral tetrads. For one of the six
words \(a\), let
\[
 \mathcal R_a
 :=\{(p,t)\in\mathcal S_+\times\mathcal S_\times:
       U_6(p,t)\text{ publishes the word }a\}
\]
the set of its one thousand eight routes, and set
\(R_9:=\mathbb Z/9\mathbb Z\). Once the oriented chart is fixed, the phase
projection has type
\begin{equation}
 q_\beta:\mathcal R_a\longrightarrow\binom{R_9}{2},
 \qquad
 q_\beta(p,t)
 =\{x-5\beta(d_t),x+5\beta(d_t)\}\subset\mathbb Z/9\mathbb Z,
 \qquad x=5(i_p-j_p),
 \label{eq:proyeccion-fano-semillas-masas-autonoma}
\end{equation}
with \((\beta(N),\beta(E),\beta(S),\beta(O))=(1,2,3,4)\). Its codomain is
the set of the thirty-six unordered pairs of \(R_9\). Each
\(X\in Q\) is a tetrad of orbits;
its route support and its cut over a fiber are defined by
\[
 \operatorname{Supp}_{\rm rutas}(X)
 :=\bigcup_{\mathcal O\in X}\mathcal O,
 \qquad
 X_b:=\operatorname{Supp}_{\rm rutas}(X)\cap q_\beta^{-1}(b),
 \qquad b\in\binom{R_9}{2}.
\]
In the multiplicative channel, the source observable is
\begin{equation}
 \rho:\mathcal S_\times\longrightarrow R_9,
 \qquad
 \rho(i,j,d)=(i+1)(j+1)\pmod9.
 \label{eq:residuo-producto-fano-semillas-rev3}
\end{equation}
If \(\pi_\times:\mathcal R_a\to\mathcal S_\times\),
\(\pi_\times(p,t)=t\), is the second projection, the observable acting
on routes is
\begin{equation}
 \rho_\times:=\rho\circ\pi_\times:\mathcal R_a\longrightarrow R_9,
 \qquad
 \rho_\times(p,t)=(i_t+1)(j_t+1)\pmod9.
 \label{eq:residuo-producto-rutas-fano-semillas-rev3}
\end{equation}
In what follows, \(\rho(X)\) abbreviates the multiset of values of
\(\rho_\times\) over \(X_b\), which is independent of the fiber \(b\). Its
profiles over the tetrads are
\[
 \rho(B_3)=3^4,
 \qquad
 \rho(B_6)=6^4,
 \qquad
 \rho(B_9)=0^4,
\]
and, for the three lateral tetrads,
\begin{equation}
 \begin{array}{c|c|c}
 r&\text{active signature of }P_r&R(P_r)\\
 \hline
 1&933&\{5,7\}\\
 2&393&\{4,8\}\\
 3&339&\{1,2\}.
 \end{array}
 \label{eq:perfiles-laterales-fano-semillas-rev3}
\end{equation}
Every residue of \(R(P_r)\) occurs twice among the four routes of
\((P_r)_b\).

\begin{theorem}[Incidence morphism of the seeds and six Fano planes]
\label{thm:semillas-seis-fano-fase-rev3}
For \(\delta\in R_9\), \(h\in\{3,6,0\}\) ---corresponding,
respectively, to \(B_3,B_6,B_9\)--- and \(r\in\{1,2,3\}\), let
\begin{equation}
 N_\delta(B_h,P_r)
 :=\sum_{b\in\binom{R_9}{2}}
 \#\left\{(x,y)\in(B_h)_b\times(P_r)_b:
 \rho_\times(y)-\rho_\times(x)=\delta\right\}.
 \label{eq:conteo-incidencia-fano-semillas-rev3}
\end{equation}
Then
\begin{equation}
 N_\delta(B_h,P_r)=
 \begin{cases}
 288,&h+\delta\in R(P_r),\\
 0,&h+\delta\notin R(P_r).
 \end{cases}
 \label{eq:conteo-288-cero-fano-semillas-rev3}
\end{equation}
If \(\delta\in R_9^\times\), for each \(B_h\) there exists a unique \(P_r\)
with count \(288\); therefore, one obtains a bijection
\[
 \theta_\delta:\{B_3,B_6,B_9\}
 \xrightarrow{\sim}\{P_1,P_2,P_3\}.
\]
Writing the image of \((B_3,B_6,B_9)\) by means of its subscripts, the six
bijections are
\begin{equation}
 \begin{array}{c|c@{\qquad}c|c}
 \delta&\theta_\delta&\delta&\theta_\delta\\
 \hline
 1&(2,1,3)&5&(2,3,1)\\
 2&(1,2,3)&7&(3,2,1)\\
 4&(1,3,2)&8&(3,1,2).
 \end{array}
 \label{eq:tabla-theta-delta-fano-semillas-rev3}
\end{equation}
Be \(P=\{P_1,P_2,P_3\}\) and, for \(p\in P\), be
\[
 \mu(p):=\bigl\{\{f,p\},P\setminus\{p\}\bigr\}
\]
the perfect matching associated over the four lateral tetrads.
Then
\begin{equation}
 \mathcal D_\delta
 :=\bigl\{\{B_3,B_6,B_9\}\bigr\}
 \cup
 \bigl\{\{B,q,q'\}:B\in\{B_3,B_6,B_9\},
       \{q,q'\}\in\mu(\theta_\delta(B))\bigr\}
 \label{eq:sistema-fano-semillas-rev3}
\end{equation}
is a system of Steiner \(\operatorname{S}(2,3,7)\). The six units
produce the six planes of Fano over these seven blocks that contain the
line transverse \(\{B_3,B_6,B_9\}\).
\end{theorem}

\begin{proof}
Fix a fiber \(b\). In \((B_h)_b\) there are four choices with residue \(h\). If \(h+\delta\in R(P_r)\), the profile \eqref{eq:perfiles-laterales-fano-semillas-rev3} provides exactly two choices in \((P_r)_b\) with that residue; if it does not belong to the support, it provides none. Since there are thirty-six fibers, \(N_\delta=36\cdot4\cdot2=288\) in the first case and zero in the second. The three lateral supports partition \(R_9^\times\), so a unit \(\delta\) determines a unique lateral support for each transversal block. Direct reading of those supports yields the table \eqref{eq:tabla-theta-delta-fano-semillas-rev3} and exhausts the \(3!=6\) bijections.

The transversal line covers the three pairs among \(B_3,B_6,B_9\). For a fixed \(B\), the two edges of the matching \(\mu(\theta_\delta(B))\) cover the four lateral points once. The three matchings are distinct and form the factorization of \(K_4\) into its three perfect matchings; hence every lateral pair occurs exactly once. The seven triples thus contain each of the twenty-one pairs of points once, which is property \(\operatorname{S}(2,3,7)\). Conversely, every Fano plane containing the transversal line assigns to each \(B\) a perfect matching distinct from the other two and recovers one of the preceding six bijections.
\end{proof}

The correspondence of sets \(\delta\mapsto\theta_\delta\) is neither an
action nor a homomorphism: \(R_9^\times\simeq C_6\) is abelian under
multiplication, whereas \(S_3\) is not, and the units do not form an additive
subgroup. The displacement \(\delta\) indexes six translations of the product
profiles.

The same calculation and the same table are verified in all six microfibers. Nevertheless,
the difference between support and role must be preserved: the support of
\(f\) coincides with one of the lateral supports and, therefore, the observable
\(\rho\) does not by itself recover the decomposition \(P\sqcup\{f\}\). That
distinction requires the additive signature already declared in the seed state.

\begin{statusbox}[title={Status and scope of Fano incidence}]
The \cref{thm:semillas-seis-fano-fase-rev3} is a finite theorem once
\(q_\beta\), the orientation, and the digital order have been fixed: it defines six exact
incidence gauges on blocks. It organizes phase and incidence in the seed domain;
it selects neither particle, mass, nor route. Its physical realization is performed
subsequently through the morphisms of the enriched ledger.
\end{statusbox}

The passage to physics preserves the chain
\begin{equation}
 \begin{aligned}
 \text{refined seed}
 &\longrightarrow(\text{signatures }+,\times;U_6;w_6;
                   \text{sheet memory})\\
 &\longrightarrow\text{CT108 record}
 \longrightarrow(\tau_{12},D_k,\Omega,P_6,\nu_{120},\nu_{270})\\
 &\longrightarrow\text{fiber and multisection}
 \longrightarrow(B_D,B_\Omega)
 \longrightarrow\text{CKM/PMNS transport}\\
 &\longrightarrow\text{character and towers}
 \longrightarrow\text{mass section}.
 \end{aligned}
 \label{eq:cadena-semilla-mezcla-masa-autonoma}
\end{equation}
CKM performs a change of basis of the quark operators,
\begin{equation}
 B_d^{(u)}=V_{\rm CKM}B_dV_{\rm CKM}^*,
 \qquad
 R_{\rm act}^{B_d^{(u)}}
 =V_{\rm CKM}R_{\rm act}^{B_d}V_{\rm CKM}^*,
 \label{eq:ckm-transporte-semillas-autonoma}
\end{equation}
and preserves spectrum and determinant. The mixing specifies the representation
of the operator in a flavour basis; the register preserves the eigenvalues and their
genealogy.

\subsection{Neutrino reader contained in the register}

The neutral sector has an integer extractor constructed exclusively
with record data. For each dodecaphase block \(E_m\), we define
\begin{equation}
 T(m)=\sum_{t\in E_m}(-1)^{\kappa(t)}\varepsilon(t)
       \mathbf1_{\{d(t)=9\}},
 \qquad
 s_C^{(m)}(\nu)=T(m)-T(m+6),
 \label{eq:lector-neutrino-ledger-autonoma}
\end{equation}
and
\begin{equation}
 \tau_m^\nu=\operatorname{sgn}s_C^{(m)}(\nu).
 \label{eq:trit-neutrino-ledger-autonoma}
\end{equation}
The four jump triangles produce
\begin{equation}
 \nu_{120}=\sum_{j=1}^{4}
 \operatorname{sgn}\!\left(\sum_{m\in\mathcal T_j}\tau_m^\nu\right),
 \qquad
 \nu_{270}=\sum_{m=1}^{12}\tau_m^\nu\tau_{m+3}^\nu.
 \label{eq:invariantes-neutrino-ledger-autonoma}
\end{equation}
Bidirectional inversion transforms
\(\nu_{120}\mapsto-\nu_{120}\) and preserves \(\nu_{270}\). This parity
defines the historical architecture of a reader independent of metrological
magnitudes used as a target. Its literal application
to the living record \(81\times108\), taking as
\(\varepsilon(t)\) the sole global sign bit preserved by that
projection, yields
\[
 s_C^{(m)}(\nu)=0,
 \qquad \tau_m^\nu=\nu_{120}=\nu_{270}=0
\]
in the eighty-one positions. Antipodal events have been
identified by projection of the register. The local result is preserved
as a recovered reader and, simultaneously, as a test of information loss: the
neutral extractor of
\cref{sec:extractor-neutro-z0-torsion-rev2} acts before that
identification and reads the local bidirectional bit, the sheet, torsion, and the
antipodal corona. The zero obtained after projection does not demote that
construction; it certifies exactly what information the global reader erases.
The frames
\(F_\ell,F_\nu\), the reader~\eqref{eq:invariantes-neutrino-ledger-autonoma},
and the aggregate form
\begin{equation}
 G_{ab}=\frac1{\sqrt{|L_a||N_b|}}
 \sum_{c\in L_a}\sum_{d\in N_b}K_{\rm registro}(c,d),
 \qquad U_{\rm int}=\operatorname{polar}(G),
 \label{eq:kernel-ledger-pmns-autonoma}
\end{equation}
Thus, the two already constructed ends of the PMNS bridge are fixed.

The cardinal equality between the fifty-six multiplicative seed vertices and the fifty-six CT108 families has undergone an initial control independent of metrological magnitudes used as targets. The natural Hadamard map on a single seed reaches only twenty-nine families; upon traversing the natural dihedral symmetries, signs, and translations, the maximum coverage is thirty-five. The intertwiner therefore acts on the enriched seed pair and preserves sheet, carry, and memory. The incidence matrix comparing shared emissions in that double domain must have rank fifty-six, be natural under orientation, and not factor through global counters. This criterion fixes the exact domain of the next construction.

\endgroup
\begingroup
\let\chapter\section
\let\section\subsection
\let\subsection\subsubsection
\let\subsubsection\paragraph
\let\clearpage\relax
\chapter{Color qutrit algebra and finite gauge connection}
\label{chap:algebra-color-qutrit}
\index{color!qutrit}\index{Weyl operators}
\index{special unitary group}\index{Wilson action}

The residual functional \(\kappa_3\) of the
\cref{chap:ocupacion-estadistica} yields three affine classes; the subsequent
atlas realizes their balance \(12+12+12\) in the
\cref{prop:color-atlas}. A triple of labels is not yet a
representation. The algebraic step requires the qutrit chart, its
orientation, and its phase character to be declared. Under those hypotheses, the representation and
its infinitesimal algebra are obtained exactly, without using masses or
chromodynamic data.

\section{Weyl pair and trace algebra \(0\)}
\label{sec:weyl-color-qutrit}

Let
\[
 V_{\rm col}=\mathbb C^3,
 \qquad
 \omega=e^{2\pi i/3},
\]
with ordered basis \((e_0,e_1,e_2)\). We define
\[
 Xe_j=e_{j+1\bmod3},
 \qquad
 Ze_j=\omega^j e_j.
\]
We have
\[
 X^3=Z^3=I,
 \qquad
 ZX=\omega XZ.
\]

\begin{theorem}[Weyl Basis and Complex Color Algebra]
\label{thm:weyl-sl3-color}
For
\[
 W_{ab}:=X^aZ^b,
 \qquad
 (a,b)\in(\mathbb Z/3\mathbb Z)^2,
\]
the nine operators \(W_{ab}\) form an orthogonal basis of
\(\operatorname{End}_{\mathbb C}(V_{\rm col})\) with respect to the
Hilbert--Schmidt product. The eight non-identity operators have zero trace and
\[
 \operatorname{span}_{\mathbb C}
 \{W_{ab}:(a,b)\ne(0,0)\}
 =\mathfrak{sl}_3(\mathbb C).
\]
\end{theorem}

\begin{proof}
The Weyl relation makes it possible to reduce any product to the form
\(X^aZ^b\). A direct calculation gives
\[
 \operatorname{tr}(W_{ab}^{\dagger}W_{cd})
 =3\,\delta_{ac}\delta_{bd}.
\]
Therefore, the nine operators are linearly independent; since
\(\dim_{\mathbb C}\operatorname{End}_{\mathbb C}(V_{\rm col})=9\), they form a
basis. The trace of \(X^aZ^b\) vanishes except for
\((a,b)=(0,0)\). The remaining eight thus form a basis of the trace-zero
subspace, whose complex dimension is eight.
\end{proof}

\begin{corollary}[Compact real form]
\label{cor:su3-color-qutrit}
The Involution
\[
 \tau(A):=-A^\dagger
\]
preserves \(\mathfrak{sl}_3(\mathbb C)\), and its fixed-point set is
\[
 \mathfrak{sl}_3(\mathbb C)^\tau
 =\{A:A^\dagger=-A,\ \operatorname{tr}A=0\}
 =\mathfrak{su}(3).
\]
This real form has dimension eight.
\end{corollary}

\begin{proof}
We have
\[
 W_{ab}^{\dagger}=\omega^{ab}W_{-a,-b}.
\]
The eight non-zero indices form four adjacent pairs. For
\[
 \mathcal R=\{(1,0),(0,1),(1,1),(1,2)\},
\]
the operators
\[
 A_u=W_u-W_u^\dagger,
 \qquad
 B_u=i(W_u+W_u^\dagger),
 \qquad u\in\mathcal R,
\]
are anti-Hermitian, traceless, and linearly independent over
\(\mathbb R\). They therefore constitute a real basis of
\(\mathfrak{su}(3)\).
\end{proof}

The residual application
\[
 \kappa_{\rm col}(r,c)
 =\mathbb C e_{\,r-c\bmod3}
\]
associates the three classes of the atlas with the three lines proper of \(Z\). The
operator \(X\) the permutes cyclically and \(Z\) records its phase. This
assignment is exact a time fixed
\[
 V_{\rm col}=\mathbb C[\mathbb F_3],
 \quad
 \text{the delta basis},
 \quad
 j\longmapsto j+1,
 \quad
 j\longmapsto\omega^j.
\]
Admissible affine changes of the chart produce unitarily conjugate
representations. The construction is therefore natural up to
unitary conjugacy; by itself, it selects neither the physical names of the
colors nor a QCD realization.

\section{Connection and Action on a Finite Cellular Complex}
\label{sec:gauge-finito-color}

\begin{construction}[Finite gauge theory relative to the complex]
\label{cons:gauge-finito-color}
Let \(\mathcal K=(V,E,F)\) be a finite oriented cellularization compatible with
the graph of routes. To each oriented edge \(e\) one assigns
\[
 U_e\in SU(3),
 \qquad
 U_{\bar e}=U_e^{-1}.
\]
A transformation \(g:V\to SU(3)\) acts by
\[
 U_e\longmapsto U_e^g
 =g_{t(e)}U_eg_{s(e)}^{-1}.
\]
For each face \(f\), the ordered product
\[
 U_f=\prod_{e\subset\partial f}U_e^{\epsilon(f,e)}
\]
defines its holonomy. With a base vertex \(v_f\) chosen, one has
\[
 U_f^g=g_{v_f}U_fg_{v_f}^{-1}.
\]
\end{construction}

\begin{proposition}[Gauge invariance and covariant transport]
\label{prop:wilson-color-finito}
For \(\beta\ge0\), the action
\[
 S_{\rm col}^{(\mathcal K,\beta)}[U]
 =\frac{\beta}{3}\sum_{f\in F}
 \left(3-\operatorname{Re}\operatorname{tr}U_f\right)
\]
is gauge-invariant and no negativa. If
\(\psi_v\in V_{\rm col}\) transforma as
\(\psi_v\mapsto g_v\psi_v\), then
\[
 \nabla_e\psi:=U_e\psi_{s(e)}-\psi_{t(e)}
\]
satisfies
\[
 \nabla_e^g\psi^g=g_{t(e)}\nabla_e\psi.
\]
\end{proposition}

\begin{proof}
The holonomy changes by conjugation and the trace is cyclic; of there the
invariance. As each \(U_f\) is unitary of dimension three,
\(\operatorname{Re}\operatorname{tr}U_f\le3\), lo that proves the not
negativity. The covariance of \(\nabla_e\) results substituting the two lawes
of transformation.
\end{proof}

\section{Positivity of Class Kernels and Harmonic Decomposition}

The preceding finite connection does not by itself imply reflection
positivity or a Yang--Mills limit. Before formulating such transfers, it is
appropriate to isolate the classical control that is closed. Let \(G\) be a compact
group and \(\widehat G\) its unitary dual. For each
\(\rho\in\widehat G\), let \(d_\rho\) be its dimension and
\(c_\rho(a)\ge0\) coefficients with sufficient decay, relative to the
growth of \(d_\rho\), for the following series to converge
absolutely.

\begin{definition}[Class nucleus]
\label{pf84:YM-07}
It is defined
\[
 K_a(g,h):=
 \sum_{\rho\in\widehat G}
 c_\rho(a)\,
 \Tr\!\bigl(\rho(g)\rho(h)^*\bigr),
 \qquad g,h\in G.
\]
The choice
\[
 c_\rho(t)=d_\rho e^{-tC_2(\rho)},
 \qquad t>0,
\]
provides the example of the heat kernel. The nonnegativity of the
coefficients does not replace the hypothesis of absolute convergence.
\end{definition}

\begin{theorem}[Positivity of Peter--Weyl]
\label{pf84:YM-08}
For \(g_1,\dots,g_n\in G\) and \(z_1,\dots,z_n\in\CC\),
\[
 \sum_{i,j=1}^n
 \overline z_i z_jK_a(g_i,g_j)\ge0.
\]
\end{theorem}

\begin{proof}
The unitarity of the representations and absolute convergence allow
the sum to be reordered:
\[
 \sum_{i,j}\overline z_i z_jK_a(g_i,g_j)
 =
 \sum_{\rho\in\widehat G}
 c_\rho(a)
 \left\lVert
 \sum_j\overline z_j\rho(g_j)
 \right\rVert_{\mathrm{HS}}^2\ge0.
\]
A negative spectral coefficient may destroy this Gram positivity;
therefore the sign forms part of the hypotheses.
\end{proof}

\begin{lemma}[Abelian Poisson Control]
\label{pf84:YM-09}
For \(G=U(1)\), \(\theta\in\RR\), and \(|a|<1\),
\[
 \sum_{n\in\ZZ}a^{|n|}e^{in\theta}
 =\frac{1-a^2}{1-2a\cos\theta+a^2}.
\]
If nonnegative Fourier coefficients are also required,
\(a\) is restricted to \(0\le a<1\).
\end{lemma}

\begin{proof}
The equality results from summing separately the two geometric series with
positive and negative indices and adding the mode \(n=0\). For \(a<0\), the
scalar identity remains valid but no longer constitutes an example of
nonnegative Peter--Weyl coefficients.
\end{proof}

\begin{remark}
The \cref{pf84:YM-07,pf84:YM-08,pf84:YM-09} are classical definitions and results
of harmonic analysis. They provide a positive control for class kernels; they
do not prove Osterwalder--Schrader positivity, a mass gap, confinement, or a
continuum limit, nor do they promote the finite gauge theory of this chapter
to quantum chromodynamics.
\end{remark}

\begin{exactbox}[title={Algebraic closure of color}]
The \cref{thm:weyl-sl3-color,cor:su3-color-qutrit} construct exactly the
chain
\[
\text{ternary residue}\longrightarrow
\text{qutrit}\longrightarrow
\mathfrak{sl}_3(\mathbb C)\longrightarrow\mathfrak{su}(3).
\]
The \cref{cons:gauge-finito-color,prop:wilson-color-finito} prolongs that
chain to a connection and a gauge action on the finite cellularization
\(\mathcal K\). This is the object proved in the chapter. The
comparison with physical chromodynamics is made subsequently and does not redefine the
qutrit, the balance \(12+12+12\), or the constructed representation.
\end{exactbox}

\endgroup
\begingroup
\let\chapter\section
\let\section\subsection
\let\subsection\subsubsection
\let\subsubsection\paragraph
\let\clearpage\relax
% Vista editorial de corpus/source/masas/02_ley_general.tex, líneas 363--491: cuerpo exacto salvo el retítulo académico del epígrafe sobre el ciclo de 108 estados. Las líneas 1--263 ya residen exactamente en 72_rev2_nucleo_ley_i.tex y las líneas 264--362 en 72_rev2_algoritmo_universal.tex.
\chapter{Systematic extension of mass theory}
\Needspace{7\baselineskip}
\section{Unified foundation of mass}
\Needspace{7\baselineskip}
\subsection{The object prior to magnitude}
The theory begins before the number. A physically individuated mass is not a real number accompanied by the name of a particle, but the terminal term of a mathematical biography. APP supplies the finite incidence space and the two correlative operations; TRIT preserves the temporal orientation of each transition; TPK transports phase, sheet, carry, boundary, route, memory, and survival. Dimensional Mechanics subsequently converts that biography into a section of the mass line. The complete order is

\[
\mathrm{APP}\longrightarrow\mathrm{TRIT}\longrightarrow\mathrm{TPK}
\longrightarrow\Gamma^{\rm surv}
\longrightarrow\gamma_{\rm obs}
\longrightarrow\mathscr L_\gamma
\longrightarrow\sigma(\gamma)
\longrightarrow\chi_{\rm mass}
\longrightarrow\mathcal T_{90,120}
\longrightarrow\mathcal L_M.
\]
Each arrow retains information used by the next. The cell does not replace the route; the route does not replace the register; a signature does not replace its genealogical register; a character does not replace the dimensional section on which it acts. This separation explains why two particles may share charge or spin and nevertheless have different masses: their enriched biographies are not the same. It also explains why particle and antiparticle share mass: the involution changes odd orientation and preserves even memory, which descends to the mass quotient.

APP consists of nine by nine published positions. Each position admits four initial orientations, so the oriented domain contains

\[
81\cdot4=324
\]
candidate routes. The internal classification produces 56 route families and 13 multisections. These cardinalities are simultaneous: 81 counts cells, 324 counts orientations, 56 counts quotients by history, and 13 counts coarse operator types. No particle table replaces these domains.

\Needspace{7\baselineskip}
\subsection{Discrete unit of the continuum that supports the mass}
The mass line is not appended to five independent theories of the continuum. It follows from the unique APP--TRIT--TPK state from which the spectral, solenoidal, gauge, coherent, and determinantal aspects jointly emerge, designated in the construction notation by \(sp,sol,gau,coh,det\). Those symbols are reading discriminants, not prior subdomains. If \(\mathcal C_{\rm disc}\) denotes the common structure, the generating arrow has the form

\[
\mathfrak E_{\rm cont}:\mathcal C_{\rm disc}
\longrightarrow
\bigl(\mathcal C_\infty;
\mathcal O_{sp},\mathcal O_{sol},\mathcal O_{gau},
\mathcal O_{coh},\mathcal O_{det}\bigr),
\]
where the five observations belong to the same limit object, share fibers, orientation, carry, memory, boundary, and readers, and are transported by the same refinement squares. Mass uses its determinant reading for the action scale, its spectral reading for the character and the towers, its gauge reading for the null and charged sectors, and its coherent and solenoidal readings for propagation and composition. None of them is introduced by presupposing that the state belongs to a classical physical class.

This unit is a typing condition: separating one of those aspects and retaining the others would produce a different object, incapable of simultaneously supporting route, action, charge, amplitude, and mass. The present theory does not reprove the complete continuum; it incorporates its unified construction as a foundation and verifies that none of the mass insertions disaggregates it.

\Needspace{7\baselineskip}
\subsection{Absolute scale and quotients}
The local determinantal reader has Smith normal form.

\[
\operatorname{SNF}(H_4)=\operatorname{diag}(1,2,2,4),
\]
from which

\[
|\operatorname{coker}H_4|=16,
\qquad
\Sigma_H=\varphi_{\rm HMT}\alpha_{\rm HMT}^{16},
\qquad
\operatorname{ord}_{10}(\Sigma_H)=-34.
\]
The decade is obtained before the SI chart. The pre and return sections of the action line are

\[
\hbar_{\rm pre}=H_5\,10^{-34}u_S,
\qquad
\hbar_{\rm ret}=\eta_{\rm ret}^{(\deg)}\,10^{-34}u_S,
\]
and its bidirectional quotient is

\[
R_{\rm act}=\frac{H_5}{\eta_{\rm ret}^{(\deg)}}.
\]
In the clock branch, we explicitly choose the return section, \(\hbar_{\rm int}:=\hbar_{\rm ret}\), as the action section. A declared clock \(t_0\) provides

\[
\omega_0=\frac{2\pi}{108t_0},
\qquad
\varepsilon_{\rm clk}=\hbar_{\rm ret}\omega_0,
\qquad
m_{\rm clk}=\varepsilon_{\rm clk}c_{\rm int}^{-2}.
\]
Adopting \(u_M^{\rm clk}:=m_{\rm clk}\) as origin of the chart is a normalization explicit. Not identifies automatically this section with a basis energetic or mass previously declared. The character produces first the operator basal

\[
\widehat M_X^{(0)}
=\mathfrak e_0m_{\rm clk}\,
\chi_{\rm full}(\sigma_X-\sigma_e,\eta_X-\eta_e),
\]
and the bidirectional reader acts afterward. For each presentation \(r\in\{D,\Omega\}\), let \(\widehat B_X^r=(\widehat B_X^r)^*\) over the finite fiber; in an infinite realization one also requires a common dense domain for functional calculus. Then

\[
\widehat M_X^{\beta,r}
=R_{\rm act}^{\widehat B_X^r/2}
\widehat M_X^{(0)}
R_{\rm act}^{\widehat B_X^r/2}.
\]
On a scalar fibre, \(\widehat B_X^r=\kappa_XI\), so that

\[
m_X^\beta=m_X^{(0)}R_{\rm act}^{\kappa_X}.
\]
If the tower coordinates are already expressed relative to the electron, \(-\eta_e\) is omitted. When a quotient is formed, the common dimensional section cancels, but the difference in bidirectional reading does not:

\[
\frac{m_Y^\beta}{m_X^\beta}
=\chi_{\rm full}(\sigma_Y-\sigma_X,\eta_Y-\eta_X)
R_{\rm act}^{\kappa_Y-\kappa_X}.
\]
Therefore, \(10^{-34}\) fixes the absolute scale of action, energy, and mass; it is not a hidden relative corrector. The fine information of a species can enter only through its character, its relative clock, or the oriented quotient of the pre- and return sections.

\Needspace{7\baselineskip}
\subsection{Internal quantity of the mass line induced by the cycle of 108 states}
The Planck--CT108 self-duality provides, in the internal normalization

\[
\hbar_{\rm int}=c_{\rm int}=t_0=1,
\]
the section

\[
m_{108}^{\rm int}
=0.058177641733144319230789692282953757\ldots .
\]
This number is a normalized quantum of the internal mass line. It is not the mass of a particle, does not by itself carry a MeV or GeV chart, and cannot be inserted as an additional species row. Its function is to verify compatibility among the action chart, the CT108 clock, and the self-duality of the mass line. The transition to a nominal mass still requires the route, signature, operator, and dimensional section of the species.

\Needspace{7\baselineskip}
\subsection{Existence of the HMT--MD output for surviving extensions and scalar-reduction criterion by fiber rank}
Given a surviving extension, a cellular reader, an integer signature, an absolutely convergent family of tower coefficients, and a dimensional mass section, a well-typed HMT--MD output exists. If the fiber has a unique fixed point compatible with the involution, the output reduces to a canonical scalar. If the fiber has higher rank, the primary output is the spectral operator on the fiber; a physical scalar mass requires the coupling functional that individuates the observed representation.

This theorem does not weaken the recovered scalar evaluations. It distinguishes three claims: the internal operator is constructed; the evaluation of a declared signature is exact; and the prospective determination of that signature is an additional theorem when it is not deduced from the reader itself. Experimental comparison is performed only after fixing which of those three claims is being published.

\Needspace{7\baselineskip}

\endgroup

\section{Universal material-realization algorithm and spectral equalizer of the readers \(K_D\) and \(K_\Omega\)}
The construction ends with the algorithm that receives only typed
genealogical data, forms the realization fiber, and requires the equalizer of
the readers before publishing a dimensional output.

\begingroup
\let\chapter\section
\let\section\subsection
\let\subsection\subsubsection
\let\subsubsection\paragraph
\let\clearpage\relax
\section{Universal Algorithm of the General Mass Law}
\Needspace{7\baselineskip}
\subsection{Cellular reader domain: surviving extension, route, record, and signature}
The entry contains only:

\[
 (\mathrm{APP},\mathrm{TRIT},\mathrm{TPK},
 \mathcal R_{324},\mathcal L_{108},
 \text{representation},\text{coupling},
 \text{common dimensional chart}).
\]
It contains no observed mass, residual mass, chosen coefficient by approximation, or external table.

\Needspace{7\baselineskip}
\subsection{Functorial construction of material realization: route fiber, intertwiner, coincidence projector, spectral operator, and scalarization}
For each physical state \(X\):

\begin{enumerate}
\item construct its descriptor of charge, generation, flavor, color, spin, statistics,
orientation, and composition;

\item form the fiber of compatible routes;
\item compute the space of intertwiners
\end{enumerate}
   \[
   \mathscr I_X=
   \operatorname{Hom}_{G_X}(V_X,\ell^2(F_X));
   \]
\begin{enumerate}
\item if \(\dim\mathscr I_X=0\), discard that identification;
\item construct the projector by Reynolds averaging and the polar factor;
\item reconstruct the genealogical registers of all routes in the support;
\item calculate signature, \(K_D\), \(K_\Omega\) and
\(\mathcal F_{\rm tower}^{2w}\);

\item apply the coupling-reading pointer, the equalizer, or the joint
spectral measure;

\item construct \(\widehat M_X\);
\item reduce to a scalar only if the variance is zero or an isolated pole exists;
\item for composites, first apply \(\mathfrak C_{12}\);
\item for widths, calculate the survival semigroup;
\item multiply by the common dimensional section;
\item freeze the output;
\item open the external comparison.
\end{enumerate}
\Needspace{7\baselineskip}
\subsection{Existence of the fiberwise realization}
The multiplicity of a representation \(\pi_X\) within a fiber is calculated by

\[
 \dim\mathscr I_X=
 \frac1{|G_X|}
 \sum_{g\in G_X}
 \overline{\chi_{\pi_X}(g)}
 |\operatorname{Fix}_{F_X}(g)|.
\]
If it is positive, an intertwiner exists. If it is greater than one, the result is a bundle of realizations, and the coupling must select a state or preserve the mixed measure. This formula replaces a nominal search by a finite algebraic condition on the 324 routes.

\Needspace{7\baselineskip}
\subsection{Spectral equalizer of the readers \(K_D\) and \(K_\Omega\) via \(Q_{\mathrm{eq}}=\mathbf1_{\{0\}}(\widehat L_D-\widehat L_\Omega)\)}
Let

\[
 L_r(\Gamma)=
 \kappa_r(\Gamma)\log R_{\rm act}
 +\sum_{h\in\mathfrak S}\eta_h^r(\Gamma)s_h,
 \qquad r\in\{D,\Omega\}.
\]
The coincidence projector is

\[
 Q_{\rm eq}=
 \mathbf1_{\{0\}}(\widehat L_D-\widehat L_\Omega).
\]
A realization admits common exact reading exactly when

\[
 P_X\le Q_{\rm eq}.
\]
If there exists an involution of the coupling that exchanges both readers, \(L_{\rm sym}\) is used. Otherwise, its joint spectral measure is published.

\Needspace{7\baselineskip}
\subsection{Cellular reader codomain: dimensionless character, tower, and mass realization}
The output may be:

\[
 \begin{cases}
 m_X\in\mathcal L_M,&\text{scalar mass},\\
 \operatorname{Spec}(\widehat M_X),&\text{multiplet},\\
 z_X=M_X-i\Gamma_X/2,&\text{resonance},\\
 0\in\mathcal L_M,&\text{zero section},\\
 (\mathcal C,F,R,\Delta E),&\text{topological sector},\\
 (L,\operatorname{Obs}_L),&\text{virtual section}.
 \end{cases}
\]
This diversity of codomains is part of the law; not an exception.

\Needspace{7\baselineskip}

\endgroup

\section{Null bifurcation and positivity}
The photonic and gluonic sectors are separated before applying the positive character.
The identity \(R_{\rm act}^{0}=1\) does not prove zero mass; nullity follows
from the gauge chart. In the matter families, the primary output is a
positive fiber operator. Only a typed physical compression produces a
scalar.

\section{Gateway to the exhaustive domain}
The definition and the theorems of the domain precede the unfolding of its
\hyperref[chap:atlas-genealogico-exhaustivo-324-rutas]{324 fichas completas}.
The cards are placed at the end of the scientific apparatus in order to preserve the
continuity of the derivation without concealing any route.
